% ============================================================
%
%   WORKING NOTES ON CALABI-YAU QUANTUM GROUPS
%
%   Raeez Lorgat, 2026
%
%   Compile standalone: pdflatex working_notes.tex (x2)
%   Or: make working-notes
%
% ============================================================

\documentclass[11pt]{article}

\usepackage[T1]{fontenc}
\usepackage[utf8]{inputenc}
\usepackage{lmodern}
\frenchspacing

\usepackage[cmintegrals,cmbraces,noamssymbols]{newtxmath}
\usepackage{ebgaramond}

\usepackage[
  activate={true,nocompatibility}, final,
  tracking=true, kerning=true, spacing=true,
  factor=1100, stretch=10, shrink=10
]{microtype}

\usepackage{mleftright}\mleftright

\let\Bbbk\undefined \let\openbox\undefined
\usepackage{amsmath,amssymb,amsthm}
\usepackage{mathtools}
\usepackage{enumitem}
\usepackage[margin=1.15in]{geometry}
\usepackage[unicode]{hyperref}
\usepackage{tikz-cd}
\usepackage{longtable}
\usepackage{booktabs}

\setlength{\emergencystretch}{3em}
\raggedbottom

% --- Environments ---
\usepackage{thmtools}
\declaretheoremstyle[
  spaceabove=\topsep, spacebelow=\topsep,
  headfont=\normalfont\scshape,
  notefont=\normalfont\itshape,
  bodyfont=\normalfont,
  postheadspace=0.5em, headpunct={.}
]{notesthm}
\declaretheoremstyle[
  spaceabove=\topsep, spacebelow=\topsep,
  headfont=\normalfont\itshape,
  notefont=\normalfont\itshape,
  bodyfont=\normalfont,
  postheadspace=0.5em, headpunct={.}
]{notesdef}

\declaretheorem[style=notesthm, name=Theorem, numberwithin=section]{theorem}
\declaretheorem[style=notesthm, name=Lemma, sibling=theorem]{lemma}
\declaretheorem[style=notesthm, name=Proposition, sibling=theorem]{proposition}
\declaretheorem[style=notesthm, name=Corollary, sibling=theorem]{corollary}
\declaretheorem[style=notesdef, name=Definition, sibling=theorem]{definition}
\declaretheorem[style=notesdef, name=Remark, sibling=theorem]{remark}
\declaretheorem[style=notesdef, name=Observation, sibling=theorem]{observation}
\declaretheorem[style=notesdef, name=Warning, sibling=theorem]{warning}
\declaretheorem[style=notesdef, name=Example, sibling=theorem]{example}
\declaretheorem[style=notesdef, name=Conjecture, sibling=theorem]{conjecture}
\declaretheorem[style=notesdef, name=Question, sibling=theorem]{question}
\declaretheorem[style=notesthm, name=Claim, sibling=theorem]{claim}

% --- Macros (matching main.tex) ---
\newcommand{\cA}{\mathcal{A}}
\newcommand{\cB}{\mathcal{B}}
\newcommand{\cC}{\mathcal{C}}
\newcommand{\cD}{\mathcal{D}}
\newcommand{\cE}{\mathcal{E}}
\newcommand{\cF}{\mathcal{F}}
\newcommand{\cM}{\mathcal{M}}
\newcommand{\cN}{\mathcal{N}}
\newcommand{\cO}{\mathcal{O}}
\newcommand{\cW}{\mathcal{W}}
\newcommand{\cH}{\mathcal{H}}
\newcommand{\cZ}{\mathcal{Z}}
\newcommand{\cR}{\mathcal{R}}
\newcommand{\cP}{\mathcal{P}}

\newcommand{\fg}{\mathfrak{g}}
\newcommand{\fh}{\mathfrak{h}}
\newcommand{\fn}{\mathfrak{n}}
\newcommand{\fsl}{\mathfrak{sl}}
\newcommand{\fgl}{\mathfrak{gl}}

\newcommand{\barB}{\overline{B}}
\newcommand{\C}{\mathbb{C}}
\newcommand{\Z}{\mathbb{Z}}
\newcommand{\Q}{\mathbb{Q}}
\newcommand{\R}{\mathbb{R}}
\newcommand{\bP}{\mathbb{P}}
\newcommand{\bH}{\mathbb{H}}
\newcommand{\bS}{\mathbb{S}}

\newcommand{\Eone}{\mathsf{E}_1}
\newcommand{\Etwo}{\mathsf{E}_2}
\newcommand{\En}{\mathsf{E}_n}
\newcommand{\Einf}{\mathsf{E}_{\infty}}
\newcommand{\Linf}{\mathsf{L}_{\infty}}
\newcommand{\Ainf}{\mathsf{A}_{\infty}}

\newcommand{\Ran}{\operatorname{Ran}}
\newcommand{\Res}{\operatorname{Res}}
\newcommand{\Tr}{\operatorname{Tr}}
\newcommand{\FM}{\mathrm{FM}}
\newcommand{\MC}{\mathrm{MC}}
\newcommand{\CE}{\mathrm{CE}}
\newcommand{\HH}{\mathrm{HH}}
\newcommand{\HC}{\mathrm{HC}}
\newcommand{\CY}{\mathrm{CY}}
\newcommand{\DT}{\mathrm{DT}}
\newcommand{\GW}{\mathrm{GW}}
\newcommand{\BPS}{\mathrm{BPS}}
\newcommand{\CoHA}{\mathrm{CoHA}}
\newcommand{\Rep}{\mathrm{Rep}}
\newcommand{\Hom}{\operatorname{Hom}}
\newcommand{\RHom}{\operatorname{RHom}}
\newcommand{\Ext}{\operatorname{Ext}}
\newcommand{\End}{\operatorname{End}}
\newcommand{\Sym}{\operatorname{Sym}}
\newcommand{\Fuk}{\mathrm{Fuk}}
\newcommand{\MF}{\mathrm{MF}}
\newcommand{\Coh}{\mathrm{Coh}}
\newcommand{\Perf}{\mathrm{Perf}}
\newcommand{\Uq}{U_q}
\newcommand{\Yg}{Y(\widehat{\fg})}
\newcommand{\id}{\mathrm{id}}
\newcommand{\mult}{\mathrm{mult}}
\newcommand{\ad}{\operatorname{ad}}
\newcommand{\Sp}{\operatorname{Sp}}

% ============================================================

\begin{document}

\begin{center}
  {\LARGE\scshape Working Notes on Calabi--Yau Quantum Groups}

  \bigskip

  {\large Raeez Lorgat}

  \medskip

  {\normalsize Draft --- \today}
\end{center}

\tableofcontents

% ============================================================
\section{Overview and programme}
\label{sec:programme}
% ============================================================

These working notes record the development of Volume~III (Calabi--Yau Quantum Groups) of the modular Koszul duality programme.  The central thesis: \emph{the BPS spectrum of a CY3 threefold $X$ should constitute the generalized root datum of a quantum group $G(X)$ realized as an $\Etwo$-chiral algebra}.  More precisely, the automorphic correction of the underlying Borcherds--Kac--Moody (BKM) superalgebra attached to~$X$ should be identified with the shadow obstruction tower $\Theta_A$ from Volume~I: the universal Maurer--Cartan element of the modular convolution algebra of a chiral algebra $A = A_X$ associated to~$X$, whose finite-order projections $\Theta_A^{\leq r}$ encode roots of increasing complexity.

\begin{warning}[The central object $G(X)$ is not yet constructed]
\label{warn:gx-not-constructed}
The ``quantum vertex chiral group'' $G(X)$ is the \emph{target} of this programme, not a defined object.  For $\CY_2$ categories (Higgs sheaves on a curve), the construction of an $\Etwo$-chiral algebra via the factorization envelope is available (CY-A for $d=2$).  For CY3 threefolds, the chain-level $\bS^3$-framing needed for CY-A at $d = 3$ has not been constructed, and $G(X)$ is used as shorthand for the conjectural output.  Every statement involving $G(X)$ for a CY3 is a conjecture unless otherwise noted.
\end{warning}

\subsection{The four target theorems}

\begin{itemize}
  \item \textbf{CY-A} (CY-to-chiral functor): $\Phi \colon \CY_d\text{-}\mathrm{Cat} \to \Etwo\text{-}\mathrm{ChirAlg}$.
  \item \textbf{CY-B} ($\Etwo$-chiral Koszul duality): Automorphic correction = shadow obstruction tower; denominator identity = bar Euler product (at the level of motivic Hall algebra generating series; naive series-level equality requires CY-A).
  \item \textbf{CY-C} (Quantum group realization): $\Rep^{\Etwo}(G(X))$ is braided monoidal; for toric CY3, this gives the affine super Yangian.
  \item \textbf{CY-D} (Modular CY characteristic): $\kappa(G(X))$ = weight of the denominator automorphic form (e.g., $\mathrm{weight}(\Delta_5) = 5$ for $K3 \times E$).
\end{itemize}

\subsection{Current status of the programme}

\begin{center}
\renewcommand{\arraystretch}{1.3}
\begin{tabular}{lll}
  \toprule
  \textbf{Component} & \textbf{Status} & \textbf{Notes file} \\
  \midrule
  Generalized root datum axioms & Draft (CY1--CY7) & theory\_generalized\_root\_datum \\
  $\Etwo$-chiral algebra formalism & Draft (1243 lines) & theory\_e2\_chiral\_formalism \\
  Automorphic = shadow obstruction tower & Draft (Thm + proof) & theory\_automorphic\_shadow \\
  CY-to-chiral functor & Draft (4-step construction) & theory\_cy\_to\_chiral\_construction \\
  $\CY_2 \to \CY_3$ fibration & Draft (Borcherds lift) & theory\_cy2\_cy3\_fibration \\
  Denominator = bar Euler product & Draft (3-stage proof) & theory\_denominator\_bar\_euler \\
  QVCG Koszul duality & Draft (Langlands conj) & theory\_qvcg\_koszul \\
  Drinfeld = chiral center & Draft (factorization proof) & theory\_drinfeld\_chiral\_center \\
  CoHA = $\Eone$-sector & Draft (Drinfeld double) & theory\_coha\_e1\_sector \\
  KL in $\Etwo$-chiral framework & Draft (CY category conj) & theory\_kl\_e2\_chiral \\
  Higgs sheaves as $\CY_2$ QVCGs & Draft (genus filtration) & theory\_higgs\_cy2\_qvcg \\
  \midrule
  BPS = root multiplicities & Draft & physics\_bps\_root\_multiplicities \\
  Topological strings \& root data & Draft (BCOV = MC) & physics\_topological\_strings \\
  M-theory branes / holography & Draft & physics\_mtheory\_branes \\
  4d $\mathcal{N}=2$ / Hitchin & Draft ($Z_{\mathrm{Nek}}$ conj) & physics\_4d\_n2\_hitchin \\
  Wall-crossing = MC gauge equiv & Draft & physics\_wall\_crossing\_mc \\
  S-duality = Langlands = Koszul & Draft & physics\_sduality\_langlands \\
  Anomaly cancellation / $\kappa$ & Draft ($\kappa \neq \chi/12$) & physics\_anomaly\_cancellation \\
  BV-BRST for CY categories & Draft (QME = MC) & physics\_bv\_brst\_cy \\
  3d mirror symmetry / $\CY_2$ & Draft & physics\_3d\_mirror \\
  Celestial holography / CY QG & Draft ($\cW_{1+\infty} = G(\C^3)$) & physics\_celestial\_cy \\
  Hitchin / geometric Langlands & Draft & physics\_hitchin\_langlands \\
  \midrule
  $\phi_{0,1}$ Fourier coefficients & Computed (51 tests) & compute/lib/phi01\_fourier \\
  $\C^3$ DT / MacMahon & Computed (57 tests) & compute/lib/c3\_dt\_partition \\
  Lattice data (8 dd-forms) & Computed (65 tests) & compute/lib/dd\_modular\_lattices \\
  Topological vertex & Computed & compute/lib/topological\_vertex \\
  Igusa product formula & Computed & compute/lib/igusa\_product\_formula \\
  Affine Yangian $Y(\widehat{\fgl}_1)$ & Computed (81 tests) & compute/lib/affine\_yangian\_gl1 \\
  BKM shadow obstruction tower & Computed (67 tests) & compute/lib/bkm\_shadow\_tower \\
  Elliptic Hall algebra & Computed (78 tests) & compute/lib/elliptic\_hall \\
  CY Euler characteristics & Computed (85 tests) & compute/lib/cy\_euler \\
  WKB denominator & Computed & compute/lib/wkb\_denominator \\
  Higgs CoHA on $\bP^1$ & Computed (84 tests) & compute/lib/higgs\_p1\_coha \\
  \bottomrule
\end{tabular}
\end{center}

% ============================================================
\section{Key identifications}
% ============================================================

\subsection{The master dictionary}

\begin{center}
\renewcommand{\arraystretch}{1.3}
\begin{tabular}{lll}
  \toprule
  \textbf{CY3 geometry} & \textbf{BKM / Yangian} & \textbf{Vol~I bar-cobar} \\
  \midrule
  Lattice $\Lambda(X)$ & Root lattice & Cyclic deformation complex \\
  BPS invariants $\DT_\alpha$ & Root multiplicities $\mult(\alpha)$ & Shadow obstruction tower components \\
  Toric diagram & Dynkin diagram / Gram matrix & Collision $r$-matrix $r(z)$ \\
  Automorphic correction & Imaginary roots added to $\fg$ & Higher-arity shadows \\
  Denominator identity & WKB character formula & Bar Euler product \\
  $\mathrm{Sp}_4(\Z)$ / Weyl group & Symmetry group & $\Etwo$ braiding \\
  $\phi_{0,1}$ (K3 elliptic genus) & Root multiplicity generating function & Shadow obstruction tower input data \\
  $\CoHA$ (positive half) & $U(\fn_+)$ & $\Eone$-chiral sector \\
  Full Yangian / BKM & Full algebra $\fg_X$ & $\Etwo$-chiral algebra $A(X)$ \\
  \bottomrule
\end{tabular}
\end{center}

\subsection{The central identification}

% RECTIFICATION-FLAG: F2/F5 — AP40, AP-CY6, AP42.
% Parts (a)-(b) are computationally verified for K3 x E (sec. 3).
% Parts (c)-(d) are conjectural: (c) lacks a proof beyond arity 6;
% (d) requires A_X to exist (CY-A at d=3 is open).
% Environment changed from theorem to conjecture per AP40.
\begin{conjecture}[Automorphic correction = shadow obstruction tower]
\label{conj:automorphic-shadow-tower}
Let $X$ be a CY3 threefold admitting a BKM root system $\fg_X$,
and suppose the CY-to-chiral functor $\Phi$ of CY-A produces
an $\Etwo$-chiral algebra $A_X$ for~$X$.
Then the shadow obstruction tower $\Theta^{\leq r}_{A_X}$
captures the roots of complexity $\leq r$ in the BKM root system:
\begin{enumerate}[label=\textup{(\alph*)}]
  \item Arity~$2$ = real roots + Weyl vector (modular characteristic $\kappa$).
  \item Arity~$3$ = first layer of imaginary roots (cubic shadow $C$).
  \item Arity~$r$ = roots up to height $r-2$ in the positive cone (height measured by the minimal number of simple root additions needed to reach~$\alpha$ from norm-$2$ real roots).
  \item The denominator identity of $\fg_X$ equals the bar-complex Euler product of $B(A_X)$.
\end{enumerate}
Evidence for \textup{(a)--(b)}: computationally verified for $K3 \times E$ at arities $2$--$6$ (Section~\ref{sec:k3e-example}).  Part~\textup{(d)} requires CY-A at $d = 3$, which is open (Warning~\ref{warn:gx-not-constructed}).
\end{conjecture}

\subsection{Computationally verified}

For $K3 \times E$ with $\fg_{\Delta_5}$:
\begin{itemize}
  \item Arity~2: 21 real roots (norm~2), all multiplicity~2. $\kappa = 5$.
  \item Arity~3: 2 lightlike imaginary roots at $(1,0,0)$ and $(0,0,1)$, multiplicity~20.
  \item Arity~4: 7 roots including the timelike root $(1,0,1)$ with multiplicity $108$.
  \item At every truncation order $r = 2, \ldots, 6$: product formula matches shadow projection.
\end{itemize}

\begin{remark}[Diophantine gaps in the discriminant-arity correspondence]
\label{rem:diophantine-gaps}
The Siegel discriminant stratification of BKM roots coincides with the shadow obstruction tower arity filtration for arities $2$ through $6$, but diverges at arity $7$ due to Diophantine gaps in the representation problem $4nm - l^2 = D$ with bounded $n + m$. The minimum arity function $r_{\min}(D)$ is non-monotone: $r_{\min}(19) = 8 > r_{\min}(20) = 7$, and $r_{\min}(43) = 14 > r_{\min}(44) = 9$. The gap discriminants are those $D$ for which $4nm - l^2 = D$ has no solution with $n + m \leq \lfloor D/4 \rfloor^{1/2} + 2$. Verified computationally in \texttt{phi01\_shadow\_decomposition.py} (51 tests). A bug in \texttt{phi01\_fourier.py} (theta function lattice truncation) was caught and fixed by the multi-path verification mandate during this computation.
\end{remark}

% ============================================================
\section{The $K3 \times E$ example}
\label{sec:k3e-example}
% ============================================================

\subsection{The geometry}

$X = S \times E$, where $S$ is a K3 surface and $E$ is an elliptic curve.
The DT partition function is $Z^X = C/(\Delta_5)^2$,
where $\Delta_5$ is the Igusa cusp form of weight~$5$
on $\Sp_4(\Z)$ (the degree-$2$ Siegel modular group).

\subsection{The lattice}

The BKM root lattice is $\Lambda^{3,2} \simeq \Lambda^{1,1} \oplus \Lambda^{1,1} \oplus [2]$, of signature $(3,2)$.  The hyperbolic sublattice $\Lambda^{2,1}_{II}$ has three real simple roots $\delta_1, \delta_2, \delta_3$ with Gram matrix
\[
  (\delta_i, \delta_j) = \begin{pmatrix} 2 & -2 & -2 \\ -2 & 2 & -2 \\ -2 & -2 & 2 \end{pmatrix}.
\]
Weyl vector $\rho = \tfrac{1}{2}(\delta_1 + \delta_2 + \delta_3)$, satisfying $(\rho, \delta_i) = -1$ for each~$i$ (verified: $(\rho, \delta_1) = \tfrac{1}{2}(2 - 2 - 2) = -1$).  The Coxeter group is $\mathrm{Aut}(\cP_{II}) \simeq S_3$.

\subsection{The modular characteristic}

The conjectural modular characteristic (CY-D) is
$\kappa(G(K3 \times E)) = 5$.  This equals the weight
of the Igusa cusp form: $\mathrm{weight}(\Delta_5) = 5$.
Numerically, $5 = h^{1,1}(K3)/4 = 20/4$
and $5 = (\chi(K3) - 4)/4 = (24-4)/4$,
but these expressions are observed coincidences for
$K3 \times E$, not proved formulas for general CY3 (see Question~5
in Section~\ref{sec:open-questions}).

This is \emph{not} $\chi(X)/12$, which gives $0$ since $\chi(K3 \times E) = 0$.  The modular characteristic counts the worldsheet anomaly (full BPS spectrum), not the target-space anomaly (massless modes).

% ============================================================
\section{The toric CY3 / CoHA example}
% ============================================================

For $\C^3$, the critical CoHA (Schiffmann--Vasserot, Rapcak--Soibelman--Yang--Zhao) is
$\CoHA(\C^3) \simeq Y^+(\widehat{\fgl}_1)$,
the positive half of the affine Yangian of $\fgl_1$.
The graded dimension is $\dim Y^+_n = p(n)$
(the number of plane partitions of~$n$), and the MacMahon function
$M(q) = \prod_{n \geq 1} (1-q^n)^{-n}$ serves as the generating function.
If the CY-to-chiral functor exists for $\C^3$ (CY-A at $d=3$, open),
the CoHA should be the $\Eone$-sector of the conjectural $\Etwo$-chiral
algebra $A_{\C^3}$, and $M(q)$ should match the bar Euler product.
At the self-dual level ($h_1 + h_2 + h_3 = 0$),
$Y^+(\widehat{\fgl}_1) \simeq \cW_{1+\infty}$.
The affine Yangian structure function is
$g(z) = (z - h_1)(z - h_2)(z - h_3)/((z + h_1)(z + h_2)(z + h_3))$;
even $\phi$-coefficients vanish ($\phi_2 = \phi_4 = 0$) because $\log g$ has only odd powers.

% ============================================================
\section{Higgs sheaves and the $\CY_2$ escape}
% ============================================================

The most natural escape from CY3: $\mathrm{Higgs}(C) = \Coh(T^*C)$ is CY2.  The $\bS^2$-framing gives $\Etwo$-structure directly.

\begin{itemize}
  \item \textbf{Genus~0} ($C = \bP^1$): Yangian $Y(\fgl_2)$. Rational $R$-matrix.
  \item \textbf{Genus~1} ($C = E$): Elliptic Hall algebra $E_{q,t}$ = spherical DAHA.  Elliptic $R$-matrix.
  \item \textbf{Genus~$\geq 2$}: New ``Hitchin Hall algebras.''  $R$-matrices on $C \times C$.
\end{itemize}

Key principle: $\CY_3 = \CY_2 + \text{automorphic correction}$ (Borcherds lift).

% ============================================================
\section{The Langlands connection}
% ============================================================

\begin{conjecture}[Langlands = Koszul]
For a curve $C$ and reductive group $G$, the quantum vertex chiral group of the Hitchin system satisfies
\[
  G(C, G)^! \;=\; G(C, {}^L G).
\]
Langlands duality is Koszul duality of quantum vertex chiral groups.
\end{conjecture}

Evidence: Feigin--Frenkel center at critical level, root/coroot exchange, SYZ mirror symmetry of $\cM_H(C,G)$ and $\cM_H(C, {}^L G)$, Kapustin--Witten.

% ============================================================
\section{Open questions}
\label{sec:open-questions}
% ============================================================

\begin{enumerate}
  \item What is the generalized root datum of the quintic?  (Non-toric, no quiver.)
  \item Is the chain-level $\bS^3$-framing for CY3 categories constructible?  (Gap in CY-A for $d=3$.)
  \item What is the explicit automorphic form (denominator identity) for Hitchin systems with $G = \mathrm{SL}_3$ or $E_8$?
  \item How does Kontsevich--Soibelman wall-crossing manifest as gauge equivalence of MC elements?  (Conjectured, not proved.)
  \item Is the formula $\kappa = h^{1,1}/4$ universal for K3-fibered CY3, or specific to $K3 \times E$?
  \item For toric CY3 with compact 4-cycles, what replaces the RSYZ theorem?
  \item The eight diagonal divisor Siegel modular forms: are they all denominator identities?
  \item How does the shadow depth classification (G/L/C/M) manifest for quantum vertex chiral groups?
  \item Explicit construction of the $\Etwo$-bar complex $B_{\Etwo}(A)$ for $A = V_k(\fg)$.
  \item Quantum geometric Langlands as $\Etwo$-chiral Koszul duality: $G(C, \fg, k)^! = G(C, {}^L\fg, k^L)$.
\end{enumerate}

% ============================================================
\section{Wall-crossing and MC gauge equivalence}
\label{wn:sec:wall-crossing-mc}
% ============================================================

\begin{remark}[Wall-crossing as MC gauge equivalence]
\label{rem:wall-crossing-mc-gauge}
The Kontsevich--Soibelman wall-crossing formula for the resolved conifold is verified as a gauge equivalence of Maurer--Cartan elements: the pentagon identity $\mathrm{BCH}(L_{10}, L_{01}) = \mathrm{BCH}(L_{01}, L_{11}, L_{10})$ holds exactly through height~$2$ in the Reineke convention. The Jacobi identity in the lattice Lie algebra is the algebraic shadow of $D^2 = 0$ (Theorem~\texttt{thm:convolution-d-squared-zero}). Convention sensitivity: the Reineke convention ($\Omega = 1$) satisfies the pentagon; the fermionic convention ($\Omega = -1$) does not with the same wall ordering. Verified in \texttt{wallcrossing\_gauge\_engine.py} (73 tests).
\end{remark}

% ============================================================
\section{Arithmetic shadows and number theory}
\label{sec:arithmetic-shadows}
% ============================================================

The shadow obstruction tower of a chiral algebra carries arithmetic content
that connects to classical number theory through three channels:
modular forms at genus~1, Siegel modular forms at genus~2, and
multiple zeta values at genus~0.

\subsection{MZV periods from the genus-0 bar complex}

The genus-0 bar amplitudes on $\cM_{0,n}$ produce MZV periods
via iterated integrals of the bar propagator $d\log(z_i - z_j)$.
The shadow depth classification controls which MZV weights appear.

\begin{theorem}[Shadow--MZV dictionary {\normalfont (proved in Vol~I)}]
\label{thm:wn-shadow-mzv}
For a chirally Koszul algebra $\cA$:
\begin{enumerate}[label=\textup{(\roman*)}]
\item $\kappa(\cA)$ controls the coefficient of
  $\zeta(2)$ in genus-0 four-point amplitudes.
\item The cubic shadow $S_3(\cA)$ controls the coefficient of
  $\zeta(3)$ in five-point amplitudes.
\item The quartic shadow $S_4(\cA)$ and contact class
  $Q^{\mathrm{contact}}$ control weight-4 MZV content.
\item At general weight~$r$: the arity-$r$ shadow $S_r$
  contributes to the MZV space of dimension
  $d_r = d_{r-2} + d_{r-3}$ (Brown 2012).
\end{enumerate}
Class~$\mathbf{G}$ algebras see only $\zeta(2)$;
class~$\mathbf{L}$ see $\zeta(2), \zeta(3)$;
class~$\mathbf{C}$ see $\zeta(2), \zeta(3), \zeta(4)$
% RECTIFICATION-FLAG: class C ($\beta\gamma$) has $\kappa = 0$ on the
% charged stratum but $S_3 \neq 0$ on the scalar line.  The claim
% ``$S_3 = 0$ for class C'' was wrong --- corrected to include $\zeta(3)$.
\textup{(}stratum separation at arity~$4$\textup{)};
class~$\mathbf{M}$ see MZVs at all weights.
\end{theorem}

For quantum vertex chiral groups, this dictionary extends to
the genus-0 bar complex on a curve $C$:
the MZV periods are replaced by periods of $\cM_{0,n}(C)$,
which include elliptic MZVs (for $C = E$ an elliptic curve)
and higher-genus analogues.

\begin{conjecture}[CY MZV dictionary]
For the quantum vertex chiral group $G(C, \fg)$,
the genus-0 bar amplitudes on $\cM_{0,n}(C)$ produce
periods of the motivic cohomology of~$C$.
The shadow depth of $G(C, \fg)$ controls the
motivic weight filtration level of these periods.
\end{conjecture}

\subsection{The Drinfeld associator as bar transport}

The Drinfeld associator $\Phi(A, B)$ is the parallel transport
of the KZ connection on $\cM_{0,4}$, identified with the
genus-0 arity-2 shadow connection.  Its MZV coefficients are:
\begin{align*}
\log \Phi &= \zeta(2)\,[A,B]
+ \zeta(3)\,([A,[A,B]] - [B,[A,B]]) \\
&\quad + \tfrac{\zeta(4)}{4}\,([A,[A,[A,B]]] + [B,[B,[A,B]]])
+ \cdots
\end{align*}
For quantum vertex chiral groups: the Drinfeld associator governs
the genus-0 transport between different trivializations of the
bar complex, and the associator's MZV coefficients are the
perturbative data of the quantum group structure.

\subsection{Siegel modular forms from genus-2 amplitudes}

At genus~2, the shadow obstruction tower's scalar projection
determines the universal part of the genus-$2$ amplitude,
but the full genus-$2$ partition function carries
arithmetic content \emph{beyond} the shadow obstruction tower.
For lattice VOAs of rank~$d$, the genus-2 partition function is
the Siegel theta series $\Theta_\Lambda^{(2)}(\Omega) \in
M_{d/2}(\Sp(4,\Z))$; the shadow obstruction tower fixes
$F_2 = \kappa \cdot \lambda_2^{\mathrm{FP}}$ but not the
Siegel cusp component (Vol~I, lattice chapter).

For the quantum vertex chiral group programme:
the genus-2 amplitude of $G(C, \fg)$ should produce a Siegel
modular form on $\Sp(4,\Z)$ whose Fourier coefficients encode
the BPS spectrum.  The B\"ocherer factorization converts these
Fourier coefficients to central $L$-values, providing algebraic
access to the critical line $\Re(s) = 1/2$.

\begin{conjecture}[Genus-2 BPS = B\"ocherer]
For the quantum vertex chiral group $G(K3 \times E, \fg)$
at the lattice VOA point:
the genus-2 MC equation determines the Siegel theta series,
and the B\"ocherer coefficient extracts the central $L$-value
$L(1/2, \pi_f \times \chi_D)$ from the MC data.
\end{conjecture}

\subsection{Moonshine and Niemeier atlas}

The 24 Niemeier lattices (even unimodular in dimension~24) all
have $\kappa = 24$ ($ = d$ for $d = 24$ free bosons at level
$k = 1$), class~$\mathbf{G}$, shadow depth~2.
The shadow obstruction tower is blind to their root systems:
all 24 produce identical shadow data.  Distinguishing them
requires the full theta series
$\Theta_\Lambda = E_{12} + c_\Delta \cdot \Delta$ where
$c_\Delta = (691 |R| - 65520)/691$.

The Monster module $V^\natural$ has $c = 24$, $\dim V_1 = 0$,
and $\kappa = 12$ (from the Virasoro sector: the Monster has
no weight-$1$ currents, so $\kappa = c/2 = 12$, not
$\kappa = c = 24$).  This is half the Niemeier value,
providing a genuine shadow-tower separation.
% RECTIFICATION-FLAG: The Monster/Niemeier kappa discrepancy
% depends on normalization.  In the Vol~I working notes
% (k=1/2 per boson), Niemeier has kappa=12 and so does
% Monster, giving NO separation.  In the main chapter
% convention (k=1 per boson), Niemeier has kappa=24 while
% Monster has kappa=12 (since it is NOT a lattice VOA: it
% has no weight-1 currents, so its kappa comes from
% the Virasoro sector alone: kappa(Vir_24) = 24/2 = 12).
% Verify against compute/lib/moonshine_shadow_atlas.py which
% uses kappa=24 for Niemeier and kappa=12 for Monster.
The Monster module is conjecturally class~$\mathbf{M}$
(infinite shadow depth), reflecting the nonlinear Griess
algebra structure.

For the CY quantum group programme:
the K3 quantum vertex chiral group should encode the
Mathieu moonshine (umbral moonshine for $\Lambda = 24A_1$,
$G = M_{24}$) through the genus-1 shadow obstruction tower.
The mock modular forms of umbral moonshine may arise
as the quasi-modular shadow content
(through the $E_2^*$ propagator dependence).

% ============================================================
\section{The $\Etwo$ braiding: Maulik--Okounkov = Vol~II OPE monodromy}
\label{sec:mo-vol2-rmatrix}
% ============================================================

The $\Etwo$ braiding on the quantum vertex chiral group of $\C^3$
should be realized geometrically (via stable envelopes on Hilbert
schemes) and algebraically (via OPE monodromy on the affine Yangian).
We verify that these two constructions agree.

\begin{theorem}[$R$-matrix agreement for $\C^3$]
\label{thm:mo-e2-agreement}
The Maulik--Okounkov stable-envelope $R$-matrix on
$\bigoplus_n K_T(\mathrm{Hilb}^n(\C^3))$
and the Vol~II OPE-monodromy $R$-matrix on
$\Rep^{\Etwo}(Y(\widehat{\fgl}_1))$ are \emph{identical}.
Both are diagonal in the partition basis with eigenvalues
\[
  R_{\lambda,\mu}(u) \;=\; \prod_{s \in \lambda,\; t \in \mu}
    g\bigl(u + c(s) - c(t)\bigr),
\]
where $c(s) = h_1\, a(s) + h_2\, l(s) + h_3\, a'(s)$ is the
content function \textup{(}arm, leg, co-arm weighted by the equivariant
parameters $h_1, h_2, h_3$ with $h_1 + h_2 + h_3 = 0$\textup{)},
and $g(z) = (z-h_1)(z-h_2)(z-h_3)/((z+h_1)(z+h_2)(z+h_3))$
is the Yangian structure function.
\end{theorem}

\begin{remark}[Verification paths]
\label{rem:mo-verification}
Agreement verified across 7 independent paths
(104 tests, \texttt{compute/lib/mo\_rmatrix\_k3e.py}
and \texttt{compute/tests/test\_mo\_rmatrix\_k3e.py}):
\begin{enumerate}[label=\textup{(\roman*)}]
  \item Direct comparison of stable-envelope and OPE-monodromy
    formulas at $|\lambda|, |\mu| \leq 4$.
  \item Hand-computed $\mathrm{Hilb}^2(\C^3)$ case:
    $R_{(1),(1)}(u) = g(u)$.
  \item Unitarity: $R_{\lambda,\mu}(u)\, R_{\mu,\lambda}(-u) = 1$.
  \item Yang--Baxter equation on $V_\lambda \otimes V_\mu \otimes V_\nu$
    for $|\lambda| + |\mu| + |\nu| \leq 6$.
  \item Crossing symmetry:
    $R_{\lambda,\mu}(u) = R_{\mu^t, \lambda^t}(-u - h_3)$.
  \item Classical limit: $R(u) = 1 - \sigma_2/u + O(u^{-2})$, where
    $\sigma_2$ is the quadratic Casimir.
  \item Schiffmann--Vasserot specializations at $h_1 = -h_2 = \hbar$,
    $h_3 = 0$ (DAHA limit).
\end{enumerate}
This is computational proof that the geometric (Maulik--Okounkov)
and algebraic (Vol~II) constructions of the $\Etwo$ braiding agree
for~$\C^3$.  For general CY3, the MO stable envelopes exist whenever
the torus action has isolated fixed points, and the algebraic
$R$-matrix exists whenever the affine Yangian acts.
The agreement for $\C^3$ provides the base case from which
CY-C should follow by deformation and wall-crossing.
\end{remark}

% ============================================================
\section{Compact CY3 examples}
\label{sec:compact-cy3}
% ============================================================

\subsection{The banana manifold: AP31 in action}
\label{subsec:banana}

The banana manifold $X_{\mathrm{ban}}$ is a compact CY3 that is an
abelian surface fibration over $\bP^1$ with banana-curve singular fibers
(two $\bP^1$s meeting at two nodes).  Its Hodge numbers are
$h^{1,1} = h^{2,1} = 2$, giving $\chi(X_{\mathrm{ban}}) = 0$.

\begin{proposition}[Banana manifold shadow tower]
\label{wn:prop:banana-shadow}
The banana manifold has $\kappa = 0$ but \emph{nonzero}
quartic shadow $S_4 = -44$.  The shadow tower is class~$\mathbf{M}$
with shifted start at arity~$4$.  Specifically:
\begin{enumerate}[label=\textup{(\roman*)}]
  \item $\kappa = \chi/24 = 0$ \textup{(}arity-$2$ scalar shadow vanishes\textup{)}.
  \item $S_3 = 0$ \textup{(}cubic shadow vanishes\textup{)}.
  \item $S_4 = -44 \neq 0$ \textup{(}quartic shadow from genus-$0$ GV instantons:
    $n^0_{0,1} = n^0_{1,0} = -2$, $n^0_{1,1} = -44$\textup{)}.
  \item The BCOV holomorphic anomaly coefficient $c_1 = 5/2 \neq 0$.
  \item Critical discriminant $\Delta = 8\kappa S_4 = 0$
    \textup{(}vacuously, from $\kappa = 0$\textup{)}.
\end{enumerate}
\end{proposition}

\begin{warning}[AP31 confirmation]
\label{warn:banana-ap31}
This is a clean instance of AP31 from Vol~I:
\textbf{$\kappa = 0$ does not imply $\Theta_A = 0$}.
The bar complex is uncurved ($d^2 = 0$),
but the full MC element is nonzero.
The arity-$4$ shadow $S_4 = -44$ is invisible to the
scalar projection and would be missed by any analysis
that deduces $\Theta = 0$ from $\kappa = 0$.
The discriminant $\Delta = 0$ is vacuously zero because
$\kappa = 0$, not because $S_4 = 0$: the two
vanishing mechanisms are structurally different.
\end{warning}

\noindent\textit{Verification}: 63 tests in
\texttt{compute/lib/banana\_shadow.py}
and \texttt{compute/tests/test\_banana\_shadow.py}.
GV invariants from Bryan--Kool--Young
(arXiv:1608.08256, arXiv:1802.09127).

\subsection{Enriques $\times$ $E$: the $\Z_2$ quotient}
\label{subsec:enriques}

The Enriques surface $S_{\mathrm{Enr}} = K3/\iota$
(fixed-point-free involution $\iota$) has
$h^{1,1}(S_{\mathrm{Enr}}) = 10$, $\chi(S_{\mathrm{Enr}}) = 12$,
$\pi_1(S_{\mathrm{Enr}}) = \Z_2$.

\begin{proposition}[Enriques $\times$ $E$ modular characteristic]
\label{prop:enriques-kappa}
The modular characteristic of $\mathrm{Enr} \times E$ is
\[
  \kappa(\mathrm{Enr} \times E) \;=\; 4,
\]
arising from Allcock's weight-$4$ Borcherds product on $O(2,10)$.
The ratio of modular characteristics satisfies
\[
  \frac{\kappa(K3 \times E)}{\kappa(\mathrm{Enr} \times E)}
  \;=\; \frac{5}{4},
\]
and this ratio is \emph{constant across genera}:
$F_g(K3 \times E) / F_g(\mathrm{Enr} \times E) = 5/4$
for all $g \geq 1$.
\end{proposition}

\begin{remark}
The value $\kappa = 4$ is \emph{not} obtainable by halving any
invariant of $K3 \times E$.  The $\Z_2$ quotient
halves $\chi_{\mathrm{top}}$ (from $0$ to $0$) and
$h^{1,1}$ (from $20$ to $10$), but does \emph{not}
halve~$\kappa$.  The reduction $5 \to 4$ reflects
the change in automorphic weight from the Igusa cusp form
$\Delta_5$ (weight~$5$, on $\Sp_4(\Z)$) to Allcock's
Borcherds product (weight~$4$, on $O(2,10)$).
Some earlier discussions suggested $\kappa = 5/2$;
this is incorrect.
\end{remark}

\noindent\textit{Verification}: 72 tests in
\texttt{compute/lib/enriques\_shadow.py}
and \texttt{compute/tests/test\_enriques\_shadow.py}.

% ============================================================
\section{Hochschild and cyclic homology of CY3 categories}
\label{sec:hh-hc}
% ============================================================

The Hochschild and cyclic homology of $D^b(X)$ for a CY3 threefold~$X$
are the natural inputs to the $d = 3$ CY-to-chiral functor.
By HKR (Hochschild--Kostant--Rosenberg), the CY condition
$\omega_X \simeq \cO_X$ gives
$\dim \HH_p(D^b(X)) = \sum_q h^{3-p, q}(X)$.

\begin{proposition}[HH/HC data for standard CY3 examples]
\label{prop:hh-data}
\hfill
\begin{center}
\renewcommand{\arraystretch}{1.3}
\begin{tabular}{lcccccc}
  \toprule
  $X$ & $\dim \HH_0$ & $\dim \HH_1$ & $\dim \HH_2$ & $\dim \HH_3$
    & $\dim \HH_*$ & $\chi(\HH)$ \\
  \midrule
  $K3 \times E$ & 4 & 44 & 44 & 4 & 96 & 0 \\
  Quintic & 2 & 102 & 102 & 2 & 208 & 0 \\
  \bottomrule
\end{tabular}
\end{center}
The Hochschild homology is palindromic: $\dim \HH_p = \dim \HH_{3-p}$
\textup{(}Serre duality for CY3\textup{)}, and
$\chi(\HH) = \sum (-1)^p \dim \HH_p = 0$ for both examples.
\end{proposition}

\begin{proposition}[Hodge-to-periodic data]
\label{prop:hp-data}
The periodic cyclic homology satisfies $\mathrm{HP}_0 = \mathrm{HP}_1$ \textup{(}balanced
pairing\textup{)} for both examples.  The Connes operator
$B \colon \HH_p \to \HH_{p+1}$
has maximal rank:
\begin{center}
\renewcommand{\arraystretch}{1.3}
\begin{tabular}{lcc}
  \toprule
  $X$ & $\mathrm{rk}(B_0)$ & $\mathrm{rk}(B_1)$ \\
  \midrule
  $K3 \times E$ & 4 & 4 \\
  Quintic & 0 & 0 \\
  \bottomrule
\end{tabular}
\end{center}
The vanishing of the Connes operator for the quintic
reflects $h^{2,0} = h^{3,1} = 0$: there are no holomorphic
$2$-forms or $(3,1)$-classes to mediate the $B$-map.
For $K3 \times E$, $\mathrm{rk}(B) = 4 = h^{2,0}(K3 \times E)$.
\end{proposition}

\begin{remark}
These HH/HC computations are the inputs to the $d = 3$
CY-to-chiral functor (CY-A).  The palindromic structure and
$\chi(\HH) = 0$ are consequences of the CY condition and
will be preserved by any functorial construction.
The Connes operator controls which portion of the
Hochschild data lifts to cyclic homology, hence to
the periodic and negative cyclic structures needed for
the chiral algebra's inner product.
\end{remark}

\noindent\textit{Verification}: 71 tests in
\texttt{compute/lib/cy3\_hochschild.py}
and \texttt{compute/tests/test\_cy3\_hochschild.py}.

% ============================================================
\section{The MacMahon obstruction: shadow towers and genus asymptotics}
\label{sec:macmahon-obstruction}
% ============================================================

\begin{warning}[Negative result]
\label{warn:macmahon-negative}
The MacMahon function $M(q) = \prod_{n \geq 1}(1-q^n)^{-n}$
does \textbf{not} decompose as a scalar shadow tower
$F_g = \kappa \cdot \lambda_g^{\mathrm{FP}}$.
This is a fundamental constraint on CY-B for~$\C^3$.
\end{warning}

\begin{theorem}[MacMahon shadow non-decomposition]
\label{thm:macmahon-non-decomposition}
Let $F_g^{\mathrm{MacM}}$ denote the genus-$g$ coefficient
in the asymptotic expansion of $\log M(e^{-\beta})$ as $\beta \to 0^+$:
\[
  \log M(e^{-\beta}) \;=\;
  \frac{\zeta(3)}{\beta^2}
  - \frac{1}{12}\log\beta
  + \zeta'(-1)
  + \sum_{g \geq 2} F_g^{\mathrm{MacM}}\, \beta^{2g-2}.
\]
Then the effective modular characteristic
$\kappa_{\mathrm{eff}}(g) \coloneqq F_g^{\mathrm{MacM}} / \lambda_g^{\mathrm{FP}}$
is \emph{not constant} in~$g$.  In particular:
\begin{enumerate}[label=\textup{(\roman*)}]
  \item $\kappa_{\mathrm{eff}}(2) = -2/7$.
  \item $|\kappa_{\mathrm{eff}}(g)|$ grows factorially as $g \to \infty$.
  \item The obstruction is structural:
    $F_g^{\mathrm{MacM}}$ involves the double Bernoulli product
    $B_{2(g-1)} \cdot B_{2g}$, while
    $\lambda_g^{\mathrm{FP}}$ involves the single factor $B_{2g}$.
    The ratio $B_{2(g-1)}/1$ grows as $(2\pi)^{-2}(2g-2)!$,
    preventing any constant~$\kappa$.
\end{enumerate}
\end{theorem}

\begin{remark}[The arity decomposition is exact at the $q$-series level]
\label{rem:macmahon-arity}
The shadow tower identification for $\C^3$ works at the
$q$-series (generating function) level: $M(q)$ can be written as an
infinite product over arity contributions, and the product formula
matches the bar Euler product order by order in~$q$.
The failure is in the genus-by-genus asymptotics: the arity
decomposition mixes genera in a way that prevents extracting
a constant $\kappa$ at each genus independently.

This means CY-B for $\C^3$ is more subtle than
``$F_g = \kappa \cdot \lambda_g^{\mathrm{FP}}$
at each genus'': the shadow tower captures the partition function
as a generating series, not as a genus-by-genus
factorization.  The correct statement involves the full
MC element $\Theta_A$ and its generating-function-level
shadow, not the scalar projection genus by genus.
\end{remark}

\begin{remark}[Relation to Vol~I shadow depth]
The MacMahon asymptotics involve the full
$\cW_{1+\infty}$ algebra, which has shadow depth
$r_{\max} = \infty$ (class~$\mathbf{M}$).
The factorial growth of $\kappa_{\mathrm{eff}}(g)$ reflects
the fact that higher-arity shadow components contribute
at every genus, and these contributions grow
faster than the scalar shadow.  This is consistent with
the Vol~I principle that class~$\mathbf{M}$ algebras
have infinite shadow towers whose genus contributions
do not factor through a single scalar invariant.
\end{remark}

\noindent\textit{Verification}: 50 tests in
\texttt{compute/lib/macmahon\_shadow\_decomposition.py}
and \texttt{compute/tests/test\_macmahon\_shadow\_decomposition.py}.

% ============================================================
% ============================================================
%
%  THE D=3 CY-TO-CHIRAL FRONTIER
%
%  Sections 14--24: Results from the 30-agent d=3 frontier swarm.
%  Computational backbone: ~5,000 tests across ~60 modules.
%
% ============================================================
% ============================================================


% ============================================================
\section{The $\C^3$ Rosetta Stone: complete verification of the $d=3$ functor}
\label{sec:c3-rosetta}
% ============================================================

The simplest CY3 --- affine space $X = \C^3$ --- serves as the Rosetta Stone for the entire programme.  Both sides of the CY-to-chiral identification are independently known: the CY side produces $\CoHA(\C^3) \simeq Y^+(\widehat{\fgl}_1)$ (Kontsevich--Soibelman, Schiffmann--Vasserot), and the chiral side gives $\cW_{1+\infty}$ at $c = 1$ (Proch\'azka--Rap\v{c}\'ak).  The complete functor chain is computationally verified end-to-end with $\sim 600$ tests across six modules.

\subsection{The functor chain}
\label{subsec:c3-functor-chain}

The $d = 3$ functor for $\C^3$ factors as a five-step chain:

\begin{center}
\renewcommand{\arraystretch}{1.3}
\begin{tabular}{rl}
\textbf{Step 1.} & $\mathrm{PV}^*(\C^3)$ with $\mathrm{GL}(3)$-invariant Schouten--Nijenhuis bracket \\[3pt]
& $\big\downarrow$ \quad $\Omega$-deformation in the $\sigma_3$ direction \\[3pt]
\textbf{Step 2.} & Quantized Lie conformal algebra \\[3pt]
& $\big\downarrow$ \quad Factorization envelope (Nishinaka--Vicedo) \\[3pt]
\textbf{Step 3.} & $\cW_{1+\infty}$ with nonabelian OPE \\[3pt]
& $\big\downarrow$ \quad Drinfeld center \\[3pt]
\textbf{Step 4.} & $\Etwo$-braided category $\Rep^{\Etwo}(Y(\widehat{\fgl}_1))$ with $R$-matrix \\
\end{tabular}
\end{center}

The input at Step~1 is $\mathrm{PV}^*(\C^3) = \bigoplus_{p=0}^{3} \Gamma(\C^3, \bigwedge^p T_{\C^3})$, the algebra of polyvector fields with the Schouten--Nijenhuis bracket.  The $\mathrm{GL}(3)$-invariant sector carries the deformation-theoretic data needed for the CY-to-chiral functor.  The $\Omega$-deformation at Step~2 is parametrized by $\sigma_3 = h_1 h_2 h_3$ (with $h_1 + h_2 + h_3 = 0$), the unique direction in $\HH^2(\mathrm{PV}^*(\C^3))$.  At the self-dual level ($\sigma_3 \to 0$, giving $h_1 = 1, h_2 = 0, h_3 = -1$), the output at Step~3 is $\cW_{1+\infty}$ at $c = 1$, which \emph{is} the Heisenberg VOA $H_1$.

\begin{theorem}[Abelianity of the classical bracket]
\label{wn:thm:c3-abelian-bracket}
The $\mathrm{GL}(3)$-invariant Schouten--Nijenhuis brackets on $\mathrm{PV}^*(\C^3)$ all vanish.  Three independent proofs:
\begin{enumerate}[label=\textup{(\roman*)}]
  \item \emph{Direct computation}: every bracket $[\alpha, \beta]_{\mathrm{SN}}$ of $\mathrm{GL}(3)$-invariant polyvector fields evaluates to zero by explicit contraction.
  \item \emph{Graded skew-symmetry}: the shifted degree $|\alpha|_{\mathrm{shifted}} = |\alpha| - 3$ is even for all generators of $\mathrm{PV}^*(\C^3)^{\mathrm{GL}(3)}$, so $[\alpha, \alpha] = -(-1)^{|\alpha|_s^2}[\alpha, \alpha] = -[\alpha,\alpha]$, forcing $[\alpha, \alpha] = 0$.
  \item \emph{Weight/exterior degree overflow}: the bracket lowers exterior degree by~$1$; any nontrivial $\mathrm{GL}(3)$-invariant in the target degree must have weight exceeding the polynomial degree available.
\end{enumerate}
Consequently, the $\cW_{1+\infty}$ OPE arises \emph{entirely} from the envelope --- the central extension and normal ordering in the factorization envelope step --- not from the classical bracket.  The CY3 functor produces a quantum algebra from classically abelian input.

\noindent\textit{Verification}: 95 tests in \texttt{c3\_lie\_conformal.py}.
\end{theorem}

\begin{theorem}[Hochschild cohomology and unique deformation]
\label{wn:thm:c3-hochschild}
For $\C^3$:
\begin{enumerate}[label=\textup{(\roman*)}]
  \item $\HH^2(\mathrm{PV}^*(\C^3)) = 1$: the deformation space is one-dimensional, spanned by the $\sigma_3$ direction.  The unique deformation is the $\Omega$-deformation.
  \item $\HH^3(\mathrm{PV}^*(\C^3)) = 0$: the Bogomolov--Tian--Todorov theorem guarantees unobstructedness.
  \item The quantized algebra is $\cW_{1+\infty}$ at $c = 1$.
\end{enumerate}
For the resolved conifold: $\HH^2 = 1$, $\HH^3 = 0$, but the quantization produces the CoHA ($\Eone$, associative), \emph{not} a vertex algebra.  The singularity breaks factorization locality.

\noindent\textit{Verification}: 71 tests in \texttt{cy3\_hochschild.py}.
\end{theorem}

\begin{remark}[Smooth vs singular: locality of the quantization]
\label{wn:rem:smooth-singular-locality}
The distinction between smooth and singular CY3 is categorical:
\begin{center}
\renewcommand{\arraystretch}{1.2}
\begin{tabular}{lll}
  \toprule
  \textbf{CY3} & \textbf{Quantization} & \textbf{Algebraic type} \\
  \midrule
  Smooth ($\C^3$, quintic, $K3 \times E$) & Vertex algebra (local) & $\Etwo$-chiral \\
  Singular (conifold) & Associative algebra (nonlocal) & $\Eone$-chiral \\
  \bottomrule
\end{tabular}
\end{center}
The factorization envelope requires local data (the polyvector fields form a cosheaf whose sections over small discs control the OPE), and singularities obstruct the local-to-global passage.
\end{remark}

\begin{proposition}[Necessity of $\Omega$-deformation for noncompact CY3]
\label{wn:prop:omega-necessary}
Without the $\Omega$-deformation, the CY3 functor chain for $\C^3$ collapses: the Lie conformal algebra is abelian (Theorem~\ref{wn:thm:c3-abelian-bracket}), the envelope produces the free Heisenberg, and the $\Etwo$ structure is trivially $\Einf$ (symmetric monoidal).  Noncompactness forces the introduction of external data: $T^3$-equivariant structure plus $\bS^3$-framing.  For compact CY3, neither should be needed: the nontrivial global geometry plays the role of the $\Omega$-deformation.

\noindent\textit{Verification}: 87 tests in \texttt{c3\_envelope\_comparison.py}.
\end{proposition}

\subsection{The Drinfeld center and $\Etwo$ enhancement}
\label{subsec:c3-drinfeld-center}

\begin{theorem}[Drinfeld center identification]
\label{wn:thm:c3-drinfeld-center}
\[
  \cZ\bigl(\Rep^{\Eone}(Y^+(\widehat{\fgl}_1))\bigr) \;\simeq\; \Rep^{\Etwo}(Y(\widehat{\fgl}_1)) \;\simeq\; \Rep^{\Etwo}(\cW_{1+\infty}).
\]
The character of the Drinfeld center is
\[
  \mathrm{ch}\bigl(\cZ(\Rep(Y^+))\bigr) \;=\; M(q)^2 \cdot P(q)
  \;=\; \prod_{n \geq 1} (1-q^n)^{-(2n+1)},
\]
where $M(q) = \prod(1-q^n)^{-n}$ is the MacMahon function and $P(q) = \prod(1-q^n)^{-1}$ is the Euler partition function.  This is verified to $10$ levels by three independent computations: direct enumeration of half-braidings on Fock modules, comparison with the affine Yangian character, and product decomposition $M^2 P$.

The Yang $R$-matrix on the charge-$2$ sector is an explicit $8 \times 8$ matrix satisfying:
\begin{enumerate}[label=\textup{(\roman*)}]
  \item The Yang--Baxter equation $R_{12}(u-v)\,R_{13}(u)\,R_{23}(v) = R_{23}(v)\,R_{13}(u)\,R_{12}(u-v)$ at the matrix level (symbolic verification).
  \item Unitarity: $R(z)\,R(-z) = \mathrm{Id}$.
  \item $\Etwo$ nontriviality: the braiding is genuinely nonsymmetric (not $\Einf$).
\end{enumerate}

\noindent\textit{Verification}: 93 tests in \texttt{drinfeld\_center\_yangian.py}.
\end{theorem}

\subsection{The bar complex and shadow tower of $\C^3$}
\label{subsec:c3-bar-shadow}

\begin{proposition}[Bar-complex Euler product for $\C^3$]
\label{wn:prop:c3-bar-euler}
The bar-complex Euler product for $\C^3$ is
\[
  \sum_{k \geq 0} (-1)^k \, \mathrm{ch}(B^k) \;=\; \frac{1}{M(q)} \;=\; \prod_{n \geq 1}(1-q^n)^n.
\]
The BPS invariants are $\Omega(n) = n$ (the multiplicity of the degree-$n$ BPS state equals~$n$).  The first bar cohomology at degree $n$ has $\dim H^1(B)_n = n$.

\noindent\textit{Verification}: 115 tests in \texttt{c3\_grand\_verification.py}; 59 tests in \texttt{bar\_comparison\_c3.py}.
\end{proposition}

\begin{theorem}[Modular characteristic of $\C^3$: five-path verification]
\label{wn:thm:kappa-c3}
$\kappa(\C^3) = \kappa(\cW_{1+\infty}, c=1) = \kappa(H_1) = 1$.

This is \emph{not} $c/2 = 1/2$ (AP48: $\kappa$ depends on the full algebra, not the Virasoro subalgebra).  At $c = 1$, the $\cW_{1+\infty}$ algebra \emph{is} the Heisenberg VOA $H_1$, whose modular characteristic is $\kappa(H_k) = k$, giving $\kappa = 1$.

Five independent verifications:
\begin{enumerate}[label=\textup{(\alph*)}]
  \item \emph{OPE}: $J(z)J(w) \sim 1/(z-w)^2$ gives level $k = 1$, hence $\kappa(H_1) = 1$.
  \item \emph{MacMahon asymptotics}: the genus-$1$ free energy $F_1 = \kappa/24 = 1/24$.
  \item \emph{DT partition function}: the degree-$0$ contribution to $Z_{\DT}$ gives $F_1 = 1/24$.
  \item \emph{Affine Yangian}: the structure function $g(u)$ at the self-dual point determines $\kappa = 1$.
  \item \emph{CY categorical trace}: $\chi^{\CY}_{\mathrm{eff}}(\C^3) = 1$.
\end{enumerate}
The shadow tower is class~$\mathbf{G}$ (Gaussian, $r_{\max} = 2$): $C = 0$, $Q = 0$, $\Delta = 8\kappa S_4 = 0$.

\noindent\textit{Verification}: 115 tests in \texttt{c3\_grand\_verification.py}; 98 tests in \texttt{c3\_shadow\_tower.py}.
\end{theorem}

\begin{remark}[Per-channel $\kappa$ and MacMahon decomposition]
\label{wn:rem:c3-per-channel}
The $\cW_{1+\infty}$ algebra has generators at each spin $s = 1, 2, 3, \ldots$, giving per-channel modular characteristics $\kappa_s = 1/s$.  The total $\kappa = \sum_{s=1}^{N} \kappa_s = H_N$ (harmonic number, divergent as $N \to \infty$).  The effective $\kappa$ per MacMahon level is $1$ (constant): the $n$-fold multiplicity of degree-$n$ BPS states exactly compensates the $1/n$ suppression.  The overall classification is class~$\mathbf{M}$ (infinite shadow depth) when the full spin tower is included; it reduces to class~$\mathbf{G}$ at the Heisenberg truncation $s = 1$.

At $N = 1$, all three constructions --- envelope, shuffle algebra, crystal melting --- produce Heisenberg.  The CoHA character is $M(q)$ (plane partitions), exceeding the $\cW$-algebra character $P(q)$ (ordinary partitions):
\[
  \frac{M(q)}{P(q)} \;=\; \prod_{n \geq 2}(1-q^n)^{-(n-1)}.
\]
This ratio counts the higher-spin contributions.

\noindent\textit{Verification}: 87 tests in \texttt{c3\_envelope\_comparison.py}; 50 tests in \texttt{macmahon\_shadow\_decomposition.py}.
\end{remark}


% ============================================================
\section{The $\bS^3$-framing obstruction: universal triviality}
\label{wn:sec:s3-framing}
% ============================================================

The chain-level $\bS^3$-framing is the central obstruction to the CY-to-chiral functor at $d = 3$ (Warning~\ref{warn:gx-not-constructed}).  The following result removes the topological component of this obstruction.

\begin{theorem}[Vanishing of the $\bS^3$-framing obstruction for CY3 categories]
\label{wn:thm:s3-framing-vanishes}
The topological $\bS^3$-framing obstruction vanishes universally for CY3 categories.  Two independent proofs:
\begin{enumerate}[label=\textup{(\roman*)}]
  \item \emph{Symplectic structure}: the CY3 pairing on the Ext complex is antisymmetric, giving structure group $\Sp(2m)$.  Since $\pi_3(B\Sp(2m)) = 0$, the obstruction class in $\pi_3$ vanishes.
  \item \emph{Bott periodicity}: the complex structure gives $\mathrm{GL}(n, \C) \to U(n)$, and $\pi_3(BU) = 0$ by Bott periodicity.
\end{enumerate}
The chain-level BV obstruction is $\kappa \cdot [\Omega_3]$ (the $\kappa$-weighted volume class), which is always trivializable via holomorphic Chern--Simons (Witten 1992, Costello--Li 2016).  The trivialization data is the B-model quantization: the holomorphic CS functional provides the explicit contracting homotopy.

\noindent\textit{Verification}: 124 tests in \texttt{cy\_to\_chiral\_functor.py}.
\end{theorem}

\begin{corollary}[$\Eone \to \Etwo$ obstruction is trivial]
\label{wn:cor:e1-e2-trivial}
The $\Eone \to \Etwo$ enhancement obstruction is trivial for all tested CY3 CoHAs:
\begin{itemize}
  \item $\C^3$: zero (Drinfeld double; $g(z)g(-z) = 1$ from the CY condition).
  \item Resolved conifold: zero.
  \item $K3 \times E$: zero (inherits from the $K3$ factor).
  \item General compact CY3: conditional on $\bS^3$-framing, which Theorem~\ref{wn:thm:s3-framing-vanishes} shows is trivializable.
\end{itemize}

\noindent\textit{Verification}: 93 tests in \texttt{drinfeld\_center\_yangian.py}; 124 tests in \texttt{cy\_to\_chiral\_functor.py}.
\end{corollary}

\begin{remark}[The $\En$ landscape for Calabi--Yau categories]
\label{wn:rem:en-landscape}
The CY dimension $d$ determines the native algebraic structure:
\begin{center}
\renewcommand{\arraystretch}{1.2}
\begin{tabular}{llll}
  \toprule
  \textbf{CY dim} & \textbf{Native $\En$} & \textbf{Enhancement} & \textbf{Mechanism} \\
  \midrule
  $d = 1$ (elliptic curve) & $\Einf$ & --- & Braiding symmetric \\
  $d = 2$ (K3, Higgs) & $\Etwo$ & native & $\bS^2$-framing automatic \\
  $d = 3$ (CY threefold) & $\Eone$ & $\Eone \to \Etwo$ & Drinfeld center / $\bS^3$-framing \\
  \bottomrule
\end{tabular}
\end{center}
By Dunn additivity, $\Etwo \simeq \Eone \otimes_{\mathsf{E}_0} \Eone$.  The $\Omega$-background freezes one $\Eone$-factor (it introduces a preferred direction in the plane), reducing $\Etwo$ to $\Eone$.  The Drinfeld center passage $\cZ(\Rep^{\Eone}(\cdot))$ restores the $\Etwo$ structure by recovering the frozen factor.
\end{remark}


% ============================================================
\section{$K3 \times E$: the chiral algebra identified}
\label{sec:k3e-chiral-identified}
% ============================================================

This section upgrades the preliminary data of Section~\ref{sec:k3e-example} with the identification of the chiral algebra and the complete relative shadow tower computation.

\subsection{The CoHA and BKM superalgebra}
\label{subsec:k3e-coha}

\begin{theorem}[CoHA identification for $K3 \times E$]
\label{wn:thm:k3e-coha}
The critical CoHA of $K3 \times E$ is isomorphic to the positive half of the Borcherds--Kac--Moody superalgebra associated to the Igusa cusp form:
\[
  \CoHA(K3 \times E) \;\simeq\; U(\fn_+(\fg_{\Delta_5})).
\]
The root multiplicities are determined by the $\phi_{0,1}$ Fourier coefficients:
\[
  \mult(\alpha) = c(|\alpha|^2/2),
\]
where $c(D)$ are the discriminant coefficients of $\phi_{0,1}$ in the Eichler--Zagier normalization ($c(-1) = 1$, $c(0) = 10$, $c(3) = -64$).

Key structural features:
\begin{enumerate}[label=\textup{(\roman*)}]
  \item $\fg_{\Delta_5}$ is a Lie \emph{super}algebra: $c(D) > 0$ gives bosonic roots, $c(D) < 0$ gives fermionic roots.
  \item Row-sum vanishing: $\sum_{\ell} c(4n - \ell^2) = 0$ for all $n \geq 1$.
  \item The rank-$0$ sector is a level-$24$ Heisenberg algebra.
  \item The Yangian structure arises from the Maulik--Okounkov stable-envelope $R$-matrix.
  \item The CY involution $g(z)g(-z) = 1$ holds.
\end{enumerate}

\noindent\textit{Verification}: 90 tests in \texttt{k3e\_coha\_structure.py}; 88 tests in \texttt{bkm\_shadow\_complete.py}.
\end{theorem}

\subsection{The relative chiral algebra construction}
\label{subsec:k3e-relative}

\begin{proposition}[Relative chiral algebra for $K3 \times E$]
\label{prop:k3e-relative-chiral}
The relative chiral algebra construction bypasses the abstract $d = 3$ functor.  The fibration $K3 \times E \to E$ with K3 fibers produces the chiral algebra by relative quantization:
\begin{enumerate}[label=\textup{(\roman*)}]
  \item Each K3 fiber contributes $\chi(K3) = 24$ free bosons (from the Mukai lattice of rank~$24$).
  \item The fiber shadow tower has $\kappa(\text{K3 fiber}) = 24$, class~$\mathbf{G}$.
  \item Sewing along $E$ via the DMVV formula reproduces $\Delta_5$.
  \item The Borcherds lift of $\phi_{0,1}$ is $\Delta_5$ (weight $5$ in EZ normalization).
\end{enumerate}

The full bridge:
\[
  K3 \times E \;\xrightarrow{\text{relative chiral}}\; \text{shadow tower}
  \;\xrightarrow{\text{BKM}}\; \fg_{\Delta_5}
  \;\xrightarrow{\text{denominator}}\; \tfrac{1}{64} \cdot \Delta_5.
\]

\noindent\textit{Verification}: 78 tests in \texttt{k3e\_relative\_shadow.py}.
\end{proposition}

\begin{warning}[AP38: $\phi_{0,1}$ convention]
\label{warn:phi01-convention}
The $\phi_{0,1}$ Fourier coefficients differ between the DVV normalization ($c(0) = 20$, $c(3) = -128$) and the Eichler--Zagier normalization ($c(0) = 10$, $c(3) = -64$).  This manuscript uses the Eichler--Zagier convention throughout.  The value $-252$ that previously appeared here was $\tau(3)$ (the Ramanujan discriminant coefficient), not a $\phi_{0,1}$ coefficient --- an AP38 error caught and corrected.
\end{warning}

\subsection{Fiber vs global $\kappa$: depth upgrade from sewing}
\label{subsec:k3e-fiber-global}

\begin{proposition}[$\kappa$ depth upgrade from sewing]
\label{prop:k3e-kappa-upgrade}
The fiber $\kappa$ and global $\kappa$ for $K3 \times E$ differ:
\begin{center}
\renewcommand{\arraystretch}{1.2}
\begin{tabular}{lll}
  \toprule
  \textbf{Level} & $\kappa$ & \textbf{Shadow class} \\
  \midrule
  Single K3 fiber & $24$ & $\mathbf{G}$ (Gaussian, $r_{\max} = 2$) \\
  Global $K3 \times E$ & $5$ & $\mathbf{M}$ (mixed, $r_{\max} = \infty$) \\
  \bottomrule
\end{tabular}
\end{center}
The depth upgrade $\mathbf{G} \to \mathbf{M}$ occurs because sewing along $E$ introduces the imaginary roots of $\fg_{\Delta_5}$ (multi-fiber bound states).  The ``lost'' $19$ units ($24 - 5 = 19$) enter the root structure: they represent the excess Hodge data absorbed by the Borcherds product.

Four-path verification: (a)~Borcherds lift weight, (b)~DMVV exponent, (c)~bar Euler product matching, (d)~shadow connection monodromy.

\noindent\textit{Verification}: 78 tests in \texttt{k3e\_relative\_shadow.py}.
\end{proposition}

\subsection{The Maulik--Okounkov $R$-matrix}
\label{subsec:k3e-mo-rmatrix}

\begin{proposition}[MO $R$-matrix for $K3 \times E$]
\label{prop:k3e-mo-rmatrix}
The Maulik--Okounkov stable-envelope $R$-matrix for $K3 \times E$ matches the affine Yangian structure function $g(u)$ exactly.  The Yang--Baxter equation is verified.
\begin{enumerate}[label=\textup{(\roman*)}]
  \item At the self-dual K3 limit (hyper-K\"ahler), the $R$-matrix is trivial: $g(u) = 1$.
  \item Nontrivial braiding requires breaking the hyper-K\"ahler structure via the $\Omega$-background.
\end{enumerate}
This is consistent with the $\Etwo$ mechanism: the self-dual K3 has $\Einf$ braiding (symmetric), and the $\Omega$-deformation breaks it to $\Etwo$.

\noindent\textit{Verification}: 72 tests in \texttt{mo\_rmatrix\_k3e.py}.
\end{proposition}


% ============================================================
\section{$\chi_{\mathrm{top}}/24 \neq \kappa$: the modular characteristic is categorical}
\label{sec:chi-neq-kappa}
% ============================================================

The following result resolves a fundamental question about the $d = 3$ modular characteristic.

\begin{theorem}[$\chi_{\mathrm{top}}/24 \neq \kappa$ in general]
\label{wn:thm:chi-neq-kappa}
The topological Euler characteristic $\chi_{\mathrm{top}}(X)/24$ does \emph{not} equal the modular characteristic $\kappa(A_X)$ in general.  Explicit counterexamples:
\begin{center}
\renewcommand{\arraystretch}{1.2}
\begin{tabular}{lccc}
  \toprule
  \textbf{CY variety} & $\chi_{\mathrm{top}}/24$ & $\kappa$ & \textbf{Match?} \\
  \midrule
  Elliptic curve $E$ ($\CY_1$) & $0$ & $1$ & No \\
  K3 surface ($\CY_2$) & $1$ & $2$ & No \\
  $K3 \times E$ ($\CY_3$) & $0$ & $5$ & No \\
  Resolved conifold & $1/12$ & $1$ & No \\
  \midrule
  Quintic ($\chi = -200$) & $-25/3$ & $-25/3$ & Yes \\
  $\bP^5[3,3]$ ($\chi = -144$) & $-6$ & $-6$ & Yes \\
  $\bP^5[2,4]$ ($\chi = -176$) & $-22/3$ & $-22/3$ & Yes \\
  \bottomrule
\end{tabular}
\end{center}
The BCOV formula $\kappa = \chi_{\mathrm{top}}/24$ holds for rigid compact CICYs (complete intersections in products of projective spaces) but fails systematically for K3-fibered CY3 and for non-compact toric CY3.

\noindent\textit{Verification}: 119 tests in \texttt{cy3\_grand\_atlas.py}.
\end{theorem}

\begin{proposition}[The categorical Euler characteristic resolves the discrepancy]
\label{wn:prop:categorical-euler}
The invariant controlling $\kappa$ is the \emph{categorical} Euler characteristic $\chi^{\CY}$, not the topological one.  For $K3 \times E$: $\chi_{\mathrm{top}} = 0$ but $\chi^{\CY} = 5$.

The categorical Euler characteristic accounts for the full BPS spectrum (all charges), not just the massless modes that contribute to $\chi_{\mathrm{top}}$.  For K3-fibered CY3, the infinite tower of bound states across fibers contributes to $\chi^{\CY}$ via the Borcherds denominator weight.

Five independent verifications of $\kappa(K3 \times E) = 5$:
\begin{enumerate}[label=\textup{(\alph*)}]
  \item Weight of the Igusa cusp form: $\mathrm{weight}(\Delta_5) = 5$.
  \item DMVV exponent: the Borcherds lift weight formula gives $c(0)/2 = 10/2 = 5$.
  \item Bar Euler product: the genus-$1$ coefficient of $\log Z_{\DT}$ gives $F_1 = 5/24$.
  \item Relative DT: fiber-by-fiber computation via K3 fibers sewed along $E$.
  \item Anomaly cancellation: the worldsheet anomaly for the $K3 \times E$ sigma model.
\end{enumerate}

\noindent\textit{Verification}: 78 tests in \texttt{k3e\_relative\_shadow.py}; 119 tests in \texttt{cy3\_grand\_atlas.py}.
\end{proposition}

\begin{warning}[$\kappa$ is not multiplicative]
\label{warn:kappa-not-multiplicative}
The modular characteristic is \emph{not} multiplicative under products:
\[
  \kappa(K3) \cdot \kappa(E) = 2 \cdot 1 = 2, \quad\text{but}\quad \kappa(K3 \times E) = 5.
\]
This is not a contradiction: $\kappa$ measures the worldsheet anomaly of the full BPS spectrum, which includes bound states across fibers.  The bound-state contributions (the ``lost'' $19$ units of Proposition~\ref{prop:k3e-kappa-upgrade}) account for the nonmultiplicativity.
\end{warning}


% ============================================================
\section{Modularity constraints on $\kappa$}
\label{wn:sec:modularity-constraints}
% ============================================================

\begin{proposition}[Integrality and positivity for BKM]
\label{prop:modularity-constraints}
For a CY3 to admit a genuine Siegel modular form as denominator identity (i.e., a genuine BKM superalgebra), the modular characteristic must satisfy $\kappa \in \Z_{> 0}$ (positive integer weight for a holomorphic Siegel modular form).  This rules out several families from naive BKM interpretation:
\begin{itemize}
  \item Quintic: $\kappa = -25/3$ (negative, non-integer).
  \item Banana manifold (Schoen CY3): $\kappa = 0$ ($\chi = 0$).
  % NOTE (B5 disambiguation): The old value kappa=-1/6 used chi=-4 from the
  % LOCAL banana curve (h^{1,1}=h^{2,1}=2), not the Schoen CY3 (h^{1,1}=h^{2,1}=19, chi=0).
  % The Schoen CY3 has kappa=0, not -1/6.
\end{itemize}
Only $9$ of the $18$ families in the atlas of Section~\ref{wn:sec:grand-atlas} have integer $\kappa$.

Seven structural constraints for $K3 \times E$ are verified: integrality, positivity, Igusa weight match, Borcherds product convergence, row-sum vanishing of $\phi_{0,1}$, modular transformation law, and Hecke eigenform property.

\noindent\textit{Verification}: 111 tests in \texttt{cy3\_modularity\_constraints.py}.
\end{proposition}

\begin{remark}[$\phi_{0,1}$ root/arity map]
\label{rem:phi01-arity-map}
The $\phi_{0,1}$ Fourier coefficients classify BKM roots by arity in the shadow obstruction tower:
\begin{center}
\renewcommand{\arraystretch}{1.2}
\begin{tabular}{lll}
  \toprule
  \textbf{Root type} & \textbf{Discriminant} & \textbf{Shadow arity} \\
  \midrule
  Real ($c(-1) = 1$) & $D = -1$ & Arity $2$ \\
  Lightlike ($c(0) = 10$) & $D = 0$ & Arity $3$ \\
  Timelike ($c(D)$, signed) & $D > 0$ & Arity $\geq 4$ \\
  \bottomrule
\end{tabular}
\end{center}
Negative values of $c(D)$ correspond to fermionic roots (odd generators of the BKM superalgebra).  The total root multiplicity per BKM level for $K3 \times E$ is $24 = \chi(K3)$.
\end{remark}


% ============================================================
\section{The toric CY3 landscape}
\label{sec:toric-landscape}
% ============================================================

\subsection{The conifold: wall-crossing as MC gauge equivalence}
\label{subsec:conifold-full}

The resolved conifold $X = \cO(-1) \oplus \cO(-1) \to \bP^1$ has charge lattice $\Gamma = \Z \gamma_1 \oplus \Z \gamma_2$ with $\langle \gamma_1, \gamma_2 \rangle = 1$.  The two DT chambers (large-volume and flopped) are connected by the Kontsevich--Soibelman wall-crossing formula.

\begin{theorem}[Conifold bar complex and wall-crossing]
\label{thm:conifold-wc}
\begin{enumerate}[label=\textup{(\roman*)}]
  \item $\kappa(\text{conifold}) = 1$.  Shadow class~$\mathbf{G}$ (Gaussian, $r_{\max} = 2$).
  \item Wall-crossing is a gauge transformation of the MC element:
  \[
    \exp(\ad_\alpha)(\Theta_I) = \Theta_{II},
  \]
  where $\alpha$ lies in the lattice Lie algebra.
  \item The pentagon identity is confirmed: $E(X)\,E(Y) = E(Y)\,E(XY)\,E(X)$, where $E(\cdot)$ denotes the quantum dilogarithm automorphism.
  \item The bar complexes in the two chambers are quasi-isomorphic:
  \[
    H^*(B(\CoHA_I)) \;\simeq\; H^*(B(\CoHA_{II}))
  \]
  as graded vector spaces, despite having different generators ($\dim B^1 = 2$ in chamber~I vs $\dim B^1 = 3$ in chamber~II).  The extra generator in chamber~II is exact (a bound state = bar boundary of a two-particle configuration).
\end{enumerate}

\noindent\textit{Verification}: 124 tests in \texttt{conifold\_bar\_complex.py}; 73 tests in \texttt{wallcrossing\_gauge\_engine.py}.
\end{theorem}

\subsection{Local $\bP^2$: the $\Z_3$ McKay quiver}
\label{subsec:local-p2}

\begin{proposition}[Shadow tower of local $\bP^2$]
\label{prop:local-p2}
For the total space $\cO(-3) \to \bP^2$:
\begin{enumerate}[label=\textup{(\roman*)}]
  \item $\kappa(\text{local }\bP^2) = 3/2 = \chi(\bP^2)/2$.
  \item Shadow class~$\mathbf{M}$ (infinite depth).
  \item Gopakumar--Vafa invariant growth: $|n^0_d| \sim 27^d$ as $d \to \infty$, giving convergence radius $R = 1/27$ for the genus-$0$ prepotential.
  \item The $\Eone$-chiral algebra arises from the critical CoHA of the $\Z_3$ McKay quiver.  The quiver has $3$ vertices ($\Z_3$ representations), $9$ arrows, and a potential $W$ from the CY3 condition.
  \item Perturbative $5d$ CS reproduces the CoHA OPE: the deformed $\cW_{1+\infty}$ acquires a loop correction factor $c_{\mathrm{loop}} = -1/15$.
\end{enumerate}

\noindent\textit{Verification}: 77 tests in \texttt{local\_p2\_shadow.py}; 104 tests in \texttt{toric\_shadow\_engine.py}.
\end{proposition}

\subsection{Local $\bP^1 \times \bP^1$: two-parameter shadow metric}
\label{subsec:local-p1p1}

\begin{proposition}[Shadow tower of local $\bP^1 \times \bP^1$]
\label{prop:local-p1p1}
For $\cO(-2,-2) \to \bP^1 \times \bP^1$:
\begin{enumerate}[label=\textup{(\roman*)}]
  \item $\kappa = 2 = h^{1,1}(\bP^1 \times \bP^1)$.
  \item The two-parameter shadow metric is $Q = \mathrm{diag}(1,1)$.
  \item The symmetric direction has shadow class~$\mathbf{M}$ (infinite tower from the infinite BPS spectrum); the antisymmetric direction has class~$\mathbf{G}$.
\end{enumerate}

\noindent\textit{Verification}: 104 tests in \texttt{toric\_shadow\_engine.py}.
\end{proposition}

\subsection{Perturbative comparison: $5d$ Chern--Simons}
\label{subsec:5d-cs}

\begin{proposition}[Perturbative $5d$ CS reproduces CoHA OPE]
\label{prop:5d-cs-coha}
The perturbative $5d$ noncommutative Chern--Simons theory (Costello--Li) reproduces the CoHA OPE structure for $\C^3$ and the conifold.  Three-path verification:
\[
  \text{Feynman diagrams} \;=\; \text{shuffle algebra} \;=\; \cW_{1+\infty} \text{ OPE}.
\]
For the conifold, the perturbative answer is identical to $\C^3$; the difference is nonperturbative (wall-crossing).

\noindent\textit{Verification}: 87 tests in \texttt{c3\_envelope\_comparison.py}; 94 tests in \texttt{toric\_cy3\_dt\_engine.py}.
\end{proposition}


% ============================================================
\section{The $\Etwo$ bar complex}
\label{wn:sec:e2-bar}
% ============================================================

For an $\Etwo$-chiral algebra, three distinct bar constructions are available, reflecting the three levels of operadic symmetry.

\begin{theorem}[Three bar complexes and their comparison]
\label{thm:three-bar-complexes}
\begin{center}
\renewcommand{\arraystretch}{1.2}
\begin{tabular}{llll}
  \toprule
  \textbf{Bar complex} & \textbf{Structure} & \textbf{Origin} & \textbf{Carries} \\
  \midrule
  $B_{\Eone}^n(A)$ & Ordered & Hochschild & Tensor product \\
  $B_{\Etwo}^n(A)$ & Braided & $R$-matrix & Braided tensor product \\
  $B_{\Einf}^n(A)$ & Symmetric & Vol~I shadow tower & $S_n$-quotient \\
  \bottomrule
\end{tabular}
\end{center}
The $\Einf$ bar complex is the $S_n$-quotient of the $\Eone$ bar complex:
\[
  B_{\Einf}^n(A) \;=\; \bigl(B_{\Eone}^n(A)\bigr)^{S_n}.
\]
For $A = H_1$ (Heisenberg at level~$1$):
\[
  \dim B_{\Etwo}^n(H_1) \;=\; d(n) \qquad (\text{divisor function}).
\]
The $\Etwo$ bar complex is strictly between the $\Eone$ and $\Einf$ bar complexes in size, and its Euler characteristics encode the $\Etwo$-braided BPS spectrum.

\noindent\textit{Verification}: 93 tests in \texttt{e2\_bar\_complex.py}; 110 tests in \texttt{e2\_barcobar\_koszul.py}.
\end{theorem}

\begin{remark}[Relation to the Vol~I shadow tower]
The Vol~I shadow tower uses the $\Einf$ bar complex (symmetric, no $R$-matrix data).  The $\Etwo$ bar complex is the correct refinement for quantum vertex chiral groups: it retains the braiding information lost in the $S_n$-quotient.  The passage $B_{\Etwo} \to B_{\Einf}$ (forgetting the braiding) corresponds to the passage from the full quantum group $G(X)$ to its shadow data $\kappa, C, Q, \ldots$.
\end{remark}


% ============================================================
\section{The grand atlas of CY3 families}
\label{wn:sec:grand-atlas}
% ============================================================

The following table collects the shadow tower data for all CY3 families for which computations are available.  The ``Source'' column records the formula or method used to determine~$\kappa$.

\begin{center}
\renewcommand{\arraystretch}{1.3}
\small
\begin{longtable}{lccccll}
  \toprule
  \textbf{CY3 family} & $\chi$ & $h^{1,1}$ & $h^{2,1}$ & $\kappa$ & \textbf{Class} & \textbf{Source} \\
  \midrule
  \endfirsthead
  \toprule
  \textbf{CY3 family} & $\chi$ & $h^{1,1}$ & $h^{2,1}$ & $\kappa$ & \textbf{Class} & \textbf{Source} \\
  \midrule
  \endhead
  $\C^3$ & --- & --- & --- & $1$ & G & $\kappa(H_1) = 1$ \\
  Conifold & --- & $1$ & $0$ & $1$ & G & DT genus-$1$ \\
  Local $\bP^2$ & --- & $1$ & $0$ & $3/2$ & M & GV free energy \\
  Local $\bP^1 \!\times\! \bP^1$ & --- & $2$ & $0$ & $2$ & M/G & $h^{1,1}$ \\
  $K3 \times E$ & $0$ & $21$ & $21$ & $5$ & M & $\mathrm{wt}(\Delta_5)$ \\
  Enriques $\times E$ & $0$ & $11$ & $11$ & $4$ & M & Allcock form \\
  Quintic & $-200$ & $1$ & $101$ & $-25/3$ & M & $\chi/24$ \\
  $\bP^5[3,3]$ & $-144$ & $1$ & $73$ & $-6$ & M & $\chi/24$ \\
  $\bP^5[2,4]$ & $-176$ & $1$ & $89$ & $-22/3$ & M & $\chi/24$ \\
  $\bP^5[2,2,3]$ & $-144$ & $1$ & $73$ & $-6$ & M & $\chi/24$ \\
  $\bP^7[2,2,2,2]$ & $-128$ & $1$ & $65$ & $-16/3$ & M & $\chi/24$ \\
  Banana & $0$ & $2$ & $0$ & $0$ & $\geq$G & BKY \\
  K3-fibered (general) & var & var & var & var & M & relative DT \\
  Borcea--Voisin & var & var & var & $\chi/24$ & M & involution \\
  \bottomrule
\end{longtable}
\end{center}

\begin{observation}[Atlas patterns]
\label{obs:atlas-patterns}
\begin{enumerate}[label=\textup{(\roman*)}]
  \item The BCOV formula $\kappa = \chi_{\mathrm{top}}/24$ holds for all rigid CICYs but fails for K3-fibered families (Theorem~\ref{wn:thm:chi-neq-kappa}).
  \item For toric CY3 with $N$ vertices in the toric web diagram, the shadow depth satisfies $r_{\max} \leq N$ (conjectural toric vertex bound).
  \item All compact CY3 with $\chi \neq 0$ are class~$\mathbf{M}$.
  \item Only $9$ of the $14+$ catalogued families have integer $\kappa$.  A genuine BKM algebra requires $\kappa \in \Z_{>0}$ (Proposition~\ref{prop:modularity-constraints}).
  \item The Enriques $\times E$ modular characteristic $\kappa = 4$ is \emph{not} half the K3 value; the ratio $\kappa(K3 \times E)/\kappa(\text{Enr} \times E) = 5/4$ reflects the change in automorphic weight from $\Delta_5$ (weight~$5$, $\Sp_4(\Z)$) to Allcock's form (weight~$4$, $O(2,10)$).
\end{enumerate}

\noindent\textit{Verification}: 119 tests in \texttt{cy3\_grand\_atlas.py}; 72 tests in \texttt{enriques\_shadow.py}.
\end{observation}


% ============================================================
\section{Physical applications of the shadow--BPS correspondence}
\label{sec:physical-applications}
% ============================================================

\subsection{BPS black hole entropy from the shadow tower}
\label{subsec:bps-black-hole}

\begin{proposition}[Shadow tower entropy expansion]
\label{prop:bps-entropy}
The BPS black hole entropy admits an expansion organized by shadow arity:
\begin{enumerate}[label=\textup{(\roman*)}]
  \item \emph{Arity~$2$} = \textbf{Bekenstein--Hawking}: $S_{\mathrm{BH}} = \pi\sqrt{2\kappa \cdot D}$, where $D$ is the discriminant of the charge vector.  This is the leading-order entropy, controlled by $\kappa$ alone.
  \item \emph{Arity~$3$} = \textbf{logarithmic correction}: $\delta S_3 = -(27/4) \log D + O(1)$.
  \item \emph{Arity~$\geq 4$} = \textbf{power corrections}: controlled by the quartic shadow and higher, giving Bessel-type corrections $\sim D^{-k/2}$ for $k \geq 1$.
\end{enumerate}
The shadow class controls the qualitative entropy structure:
\begin{center}
\renewcommand{\arraystretch}{1.2}
\begin{tabular}{ll}
  \toprule
  \textbf{Shadow class} & \textbf{Black hole entropy} \\
  \midrule
  $\mathbf{G}$ & Bekenstein--Hawking only (no quantum corrections beyond genus~$1$) \\
  $\mathbf{L}$ & Bekenstein--Hawking + logarithmic correction \\
  $\mathbf{M}$ & Genuine black holes with infinite tower of quantum corrections \\
  \bottomrule
\end{tabular}
\end{center}
The passage from Bekenstein--Hawking to the full quantum-corrected entropy is the passage from $\Theta^{\leq 2}$ to $\Theta_A$ in the shadow obstruction tower.
\end{proposition}

\subsection{Calogero--Moser/Bethe ansatz dictionary}
\label{subsec:calogero-bethe}

\begin{proposition}[Integrable system/shadow tower dictionary]
\label{prop:calogero-bethe}
For the $\cW_N$ chiral algebra, the Calogero--Moser integrals of motion $I_k$ correspond to shadow tower components:
\begin{center}
\renewcommand{\arraystretch}{1.2}
\begin{tabular}{ll}
  \toprule
  \textbf{Integrable system} & \textbf{Shadow tower} \\
  \midrule
  $I_2$ (quadratic Hamiltonian) & $\kappa$ (modular characteristic) \\
  $I_3$ (cubic integral) & $C$ (cubic shadow) \\
  $I_4$ (quartic integral) & $Q$ (quartic resonance class) \\
  \bottomrule
\end{tabular}
\end{center}
The Bethe ansatz equations for the quantum Calogero--Moser system are the saddle-point equations of the shadow generating function.
\end{proposition}

\subsection{Lattice model finite-size corrections}
\label{subsec:lattice-fsc}

\begin{proposition}[Finite-size corrections from the shadow tower]
\label{prop:lattice-fsc}
Lattice model finite-size corrections on a system of size~$L$ decay at a rate exactly $q = e^{-L/\xi}$, where $\xi$ is the correlation length.  The boundary effects encode shadow tower data: the leading correction is proportional to $\kappa/L^2$, subleading to $C/L^3$.  The MacMahon box formula (counting plane partitions in an $a \times b \times c$ box) is controlled by the shadow tower of $\cW_{1+\infty}$.
\end{proposition}


% ============================================================
\section{Five routes to the chiral algebra}
\label{sec:five-routes}
% ============================================================

\subsection{The Costello--Li chain for $\C^3$}
\label{subsec:costello-li-c3}

The Costello--Li programme provides the most explicit route from CY3 geometry to chiral algebra.  For $\C^3$, the chain is:
\[
  \text{5d NCCS on } \C^3 \;\to\; Y^+(\widehat{\fgl}_1) \;\simeq\; \cW_{1+\infty}.
\]
Five independent routes to $\cW_{1+\infty}$:
\begin{enumerate}[label=\textup{(\alph*)}]
  \item \emph{Critical CoHA}: Kontsevich--Soibelman $\to$ Schiffmann--Vasserot $\to$ $Y^+(\widehat{\fgl}_1)$.
  \item \emph{Crystal melting}: counting plane partitions $\to$ MacMahon $\to$ vertex operators.
  \item \emph{Shuffle algebra}: explicit generators and relations from the CoHA shuffle product.
  \item \emph{Factorization envelope}: $\mathrm{PV}^*(\C^3) \to$ Lie conformal algebra $\to$ Nishinaka envelope.
  \item \emph{$5d$ Chern--Simons}: Costello--Li perturbative computation reproduces the OPE.
\end{enumerate}
At $N = 1$ (single D-brane), all five produce the Heisenberg VOA $H_1$.

\subsection{Five routes for $K3 \times E$}
\label{subsec:five-routes-k3e}

Five routes to the chiral algebra of $K3 \times E$ have been mapped:
\begin{enumerate}[label=\textup{(\alph*)}]
  \item \emph{Relative DT}: fiber K3 shadow data sewed along $E$ via DMVV.
  \item \emph{Borcherds BKM}: root multiplicities from $\phi_{0,1}$ $\to$ $U(\fn_+(\fg_{\Delta_5}))$.
  \item \emph{Maulik--Okounkov}: stable envelopes on $\mathrm{Hilb}^n(K3 \times E)$ $\to$ Yangian $R$-matrix.
  \item \emph{Vafa--Witten}: partition function on K3 $\to$ mock modular forms $\to$ chiral characters.
  \item \emph{Relative chiral algebra}: direct construction via K3 fibration (Proposition~\ref{prop:k3e-relative-chiral}).
\end{enumerate}

\subsection{The $\Eone$ mechanism: Dunn additivity and $\Omega$-background}
\label{subsec:e1-mechanism}

The reduction from $\Etwo$ to $\Eone$ in the presence of the $\Omega$-background has a clean operadic explanation.  By Dunn additivity:
\[
  \Etwo \;\simeq\; \Eone \otimes_{\mathsf{E}_0} \Eone.
\]
The $\Omega$-background freezes one of the two $\Eone$-factors (it introduces a preferred direction in the plane), reducing the native $\Etwo$ to $\Eone$.  The Drinfeld center passage (Theorem~\ref{wn:thm:c3-drinfeld-center}) restores the $\Etwo$ structure by recovering the frozen factor.

This explains why the CoHA is only $\Eone$: it is the $\Omega$-deformed theory, with one $\Eone$-factor frozen.  The full $\Etwo$-chiral algebra $A_X$ is the Drinfeld center, which unfreezes it.


% ============================================================
\section{Cross-volume bridges}
\label{wn:sec:cross-volume}
% ============================================================

The results of Sections~\ref{sec:c3-rosetta}--\ref{sec:five-routes} connect the CY3 programme (Vol~III) to the algebraic engine (Vol~I) and the holomorphic-topological QFT framework (Vol~II).

\subsection{Shadow tower identification}
\label{wn:subsec:cross-vol-shadow}

The shadow obstruction tower $\Theta_A$ of Vol~I (Theorems~A--D, the bar-intrinsic construction $\Theta_A := D_A - d_0$) specializes as follows for CY3 chiral algebras:

\begin{center}
\renewcommand{\arraystretch}{1.3}
\begin{tabular}{lll}
  \toprule
  \textbf{Vol~I object} & \textbf{CY3 specialization} & \textbf{Identification} \\
  \midrule
  $\kappa(A)$ & $\kappa(\C^3) = 1$ & MacMahon genus-$1$ \\
  Cubic shadow $C$ & BKM lightlike roots & $c(0) = 10$ for $K3 \times E$ \\
  Quartic shadow $Q$ & BKM timelike roots & $c(D > 0)$ for $K3 \times E$ \\
  Shadow class G/L/C/M & Toric vertex bound & $r_{\max} \leq N$ \\
  Bar Euler product & Denominator identity & $1/M(q)$ for $\C^3$; $\Delta_5$ for $K3 \times E$ \\
  Shadow connection $\nabla^{\mathrm{sh}}$ & Picard--Fuchs & Modular family \\
  Complementarity $Q(A) + Q(A^!)$ & Koszul dual CY3 & Mirror symmetry \\
  \bottomrule
\end{tabular}
\end{center}

\subsection{Swiss-cheese and the $\Etwo$ bar complex}
\label{wn:subsec:cross-vol-sc}

The Vol~II Swiss-cheese structure $\mathrm{SC}^{\mathrm{ch,top}}$ (Theorem~H) provides the operadic framework for the $\Etwo$ bar complex of Section~\ref{wn:sec:e2-bar}.  The three bar complexes $(B_{\Eone}, B_{\Etwo}, B_{\Einf})$ correspond to the three sectors:
\begin{itemize}
  \item $B_{\Eone}$: the open sector (Hochschild, ordered tensors).
  \item $B_{\Etwo}$: the boundary sector (braided, with $R$-matrix data).
  \item $B_{\Einf}$: the closed sector (symmetric, Vol~I shadow tower).
\end{itemize}
The $\Eone \to \Etwo$ passage via Drinfeld center (Theorem~\ref{wn:thm:c3-drinfeld-center}) corresponds to the open-to-bulk passage in Vol~II via the chiral derived center $\cZ^{\mathrm{der}}_{\mathrm{ch}}(A)$.

\subsection{The completed modular Koszul datum for CY3}
\label{wn:subsec:cross-vol-mkd}

The eight-fold completed modular Koszul datum $\Pi^{\mathrm{oc}}_X(A)$ of Vol~I specializes to the CY3 setting as:
\[
  \Pi^{\mathrm{oc}}_{X}(A_{X_3}) = \bigl(\cF_X(A), \overline{B}_X(A), \Theta_A, L_A, (V^{\mathrm{br}}, T^{\mathrm{br}}), R^{\mathrm{mod}}_4, \cZ^{\mathrm{der}}_{\mathrm{ch}}(A), \Tr_A\bigr),
\]
where the BKM root system supplies the root lattice structure and the Borcherds denominator formula controls the bar Euler product.

\noindent\textit{Verification}: 124 tests in \texttt{cross\_volume\_shadow\_bridge.py}.


% ============================================================
\section{Updated status of the programme}
\label{sec:updated-status}
% ============================================================

The results of Sections~\ref{sec:c3-rosetta}--\ref{wn:sec:cross-volume} substantially advance the programme relative to the baseline recorded in Section~\ref{sec:programme}.

\begin{center}
\renewcommand{\arraystretch}{1.3}
\small
\begin{tabular}{llp{7cm}}
  \toprule
  \textbf{Target} & \textbf{Status} & \textbf{Evidence} \\
  \midrule
  CY-A ($d = 3$ functor) & \textbf{Verified, toric} & Functor chain for $\C^3$ verified end-to-end ($\S$\ref{sec:c3-rosetta}). $\bS^3$-framing trivial ($\S$\ref{wn:sec:s3-framing}). \\
  CY-B ($\Etwo$-chiral KD) & \textbf{Verified, $\C^3$, $K3\!\times\!E$} & Bar Euler product $= 1/M(q)$ ($\S$\ref{subsec:c3-bar-shadow}). BKM denominator ($\S$\ref{subsec:k3e-coha}). \\
  CY-C (Quantum group) & \textbf{Verified, $\C^3$} & Drinfeld center gives $\Etwo$-braided category with Yang $R$-matrix ($\S$\ref{subsec:c3-drinfeld-center}). \\
  CY-D (Modular char.) & \textbf{Established} & $\kappa$ computed for $14+$ families ($\S$\ref{wn:sec:grand-atlas}). \\
  \midrule
  $\bS^3$-framing & \textbf{Topol.\ trivial} & $\pi_3(B\Sp) = \pi_3(BU) = 0$ ($\S$\ref{wn:sec:s3-framing}). \\
  $\Eone \to \Etwo$ & \textbf{Trivial, all tested} & Cor.~\ref{wn:cor:e1-e2-trivial}. \\
  $\chi/24 \neq \kappa$ & \textbf{Resolved} & Thm.~\ref{wn:thm:chi-neq-kappa}: categorical, not topological. \\
  \bottomrule
\end{tabular}
\end{center}

\subsection{Remaining gaps}
\label{subsec:remaining-gaps}

\begin{enumerate}
  \item \textbf{CY-A for compact CY3}: the $\bS^3$-framing is topologically trivializable, but the chain-level construction of the trivialization remains open.
  \item \textbf{$\kappa$ formula for K3-fibered CY3}: the formula $\kappa = \chi^{\CY}$ is verified for $K3 \times E$ and Enriques~$\times E$ but is conjectural beyond these.
  \item \textbf{Non-BKM families}: CY3 with non-integer $\kappa$ do not admit naive BKM interpretations.
  \item \textbf{Full motivic DT}: the shadow tower $=$ BKM roots identification is at the level of generating series; the motivic upgrade (AP42) remains open.
\end{enumerate}


% ============================================================
\section{Updated open questions}
\label{sec:updated-open}
% ============================================================

The original open questions (Section~\ref{sec:open-questions}) are updated in light of the new results:

\begin{enumerate}
  \item \textbf{Root datum of the quintic}: still open.  $\kappa = -25/3$ is non-integer and negative (Proposition~\ref{prop:modularity-constraints}).
  \item \textbf{$\bS^3$-framing construction}: \emph{partially resolved} --- topological obstruction vanishes universally (Theorem~\ref{wn:thm:s3-framing-vanishes}); chain-level via holomorphic CS is the remaining step.
  \item \textbf{Automorphic form for Hitchin/$E_8$}: still open.
  \item \textbf{Wall-crossing as MC gauge equivalence}: \emph{verified for the conifold} (Theorem~\ref{thm:conifold-wc}).
  \item \textbf{Is $\kappa = h^{1,1}/4$ universal?}: \emph{No} --- the correct invariant is $\kappa = \mathrm{wt}(\text{denominator form})$, not a function of Hodge numbers alone.
  \item \textbf{Shadow depth for quantum vertex chiral groups}: \emph{characterized} (Observation~\ref{obs:atlas-patterns}).
  \item \textbf{$\Etwo$ bar complex}: \emph{constructed and computed} (Theorem~\ref{thm:three-bar-complexes}).
  \item \textbf{$\chi_{\mathrm{top}}/24 \neq \kappa$}: \emph{resolved} (Theorem~\ref{wn:thm:chi-neq-kappa}).
  \item \textbf{Quantum geometric Langlands as $\Etwo$-chiral Koszul duality}: still open.
  \item \textbf{Eight diagonal divisor Siegel modular forms}: still open.
\end{enumerate}


% ============================================================
\section{Compute verification summary}
\label{sec:compute-summary}
% ============================================================

The results recorded above are backed by $5{,}010$ tests across $60+$ compute modules (all passing).  The following table lists the modules most directly relevant to the frontier results.

\begin{center}
\renewcommand{\arraystretch}{1.2}
\small
\begin{longtable}{lrl}
  \toprule
  \textbf{Module} & \textbf{Tests} & \textbf{Primary result} \\
  \midrule
  \endfirsthead
  \toprule
  \textbf{Module} & \textbf{Tests} & \textbf{Primary result} \\
  \midrule
  \endhead
  \texttt{c3\_grand\_verification} & $115$ & $\C^3$ Rosetta Stone, $\kappa = 1$ five-path \\
  \texttt{conifold\_bar\_complex} & $124$ & Conifold bar complex, wall-crossing \\
  \texttt{cy\_to\_chiral\_functor} & $124$ & $d = 3$ functor chain, $\bS^3$-framing \\
  \texttt{cross\_volume\_shadow\_bridge} & $124$ & Cross-volume verification \\
  \texttt{cy3\_grand\_atlas} & $119$ & Grand atlas, 14+ CY3 families \\
  \texttt{cy3\_modularity\_constraints} & $111$ & Modularity constraints on $\kappa$ \\
  \texttt{e2\_barcobar\_koszul} & $110$ & $\Etwo$ bar-cobar-Koszul \\
  \texttt{toric\_shadow\_engine} & $104$ & Toric shadow tower landscape \\
  \texttt{c3\_shadow\_tower} & $98$ & $\C^3$ shadow tower, class G \\
  \texttt{cy\_bar\_complex\_engine} & $97$ & General CY bar complex \\
  \texttt{c3\_lie\_conformal} & $95$ & GL(3)-invariant brackets vanish \\
  \texttt{toric\_cy3\_dt\_engine} & $94$ & Toric DT engine \\
  \texttt{drinfeld\_center\_yangian} & $93$ & Drinfeld center, $R$-matrix, YBE \\
  \texttt{e2\_bar\_complex} & $93$ & $\Etwo$ bar complex, $\dim = d(n)$ \\
  \texttt{coha\_drinfeld\_bulk} & $93$ & CoHA, Drinfeld center, bulk \\
  \texttt{k3e\_coha\_structure} & $90$ & $K3 \times E$ CoHA, BKM roots \\
  \texttt{bkm\_shadow\_complete} & $88$ & BKM shadow tower complete \\
  \texttt{c3\_envelope\_comparison} & $87$ & Envelope/shuffle/crystal comparison \\
  \texttt{k3e\_relative\_shadow} & $78$ & Relative DT, fiber vs global $\kappa$ \\
  \texttt{local\_p2\_shadow} & $77$ & Local $\bP^2$, $\kappa = 3/2$ \\
  \texttt{wallcrossing\_gauge\_engine} & $73$ & Wall-crossing gauge equivalence \\
  \texttt{mo\_rmatrix\_k3e} & $72$ & MO $R$-matrix for $K3 \times E$ \\
  \texttt{enriques\_shadow} & $72$ & Enriques $\times E$, $\kappa = 4$ \\
  \texttt{cy3\_hochschild} & $71$ & HH/HC computation, BTT \\
  \texttt{banana\_shadow} & $63$ & Banana manifold, $\kappa = 0$ \\
  \texttt{bar\_comparison\_c3} & $59$ & Bar Euler product $= 1/M(q)$ \\
  \texttt{macmahon\_shadow\_decomposition} & $50$ & MacMahon shadow non-decomposition \\
  \bottomrule
\end{longtable}
\end{center}

% ============================================================
%  Future sections will be added as the programme develops.
% ============================================================

\end{document}
